
%%FILE blueprint/src/addition.tex
\chapter{Addition, negation, and surreal numbers}

\section{Addition and Negation on games}

\begin{definition}[Recursive sum]
  \label{def:game-add}
  \lean{Game.add,add_left,add_right,mem_add_left_iff,mem_add_right_iff}
  \leanok
  \uses{def:game,def:birthday,lem:birthday-options}
  For games \(x=\{X^L\mid X^R\}\) and \(y=\{Y^L\mid Y^R\}\), define
  \[
    x+y=
    \bigl\{\,x^L+y,\ x+y^L\ \bigm|\ x^R+y,\ x+y^R\,\bigr\},
  \]
  with one option for every member of the indicated finite lists.
\end{definition}

\begin{theorem}[Basic addition laws]
  \label{thm:game-add-laws}
  \lean{Game.add_zero',Game.zero_add',Game.add_zero,Game.zero_add,Game.add_comm,Game.add_assoc}
  \leanok
  For all games \(a,b,c\),
  \[
    a+0=a,\qquad 0+a=a,\qquad a+b\GameEq b+a,
    \qquad (a+b)+c=a+(b+c).
  \]
  The zero and associativity identities hold as literal equalities of game trees; commutativity
  is asserted up to game equivalence.
\end{theorem}
\begin{proof}
  \leanok
  \uses{def:game,def:game-equivalence,lem:birthday-options,def:game-inductions,thm:game-order-calculus,thm:option-extensionality,def:game-add}
  For \(a+0=a\), use the one-game induction scheme of
  Definition~\ref{def:game-inductions}.  Write \(a=\{A^L\mid A^R\}\).  By
  Definition~\ref{def:game-add}, the left and right option lists of \(a+0\)
  are respectively
  \[
    [\,u+0:u\in A^L\,],
    \qquad
    [\,v+0:v\in A^R\,].
  \]
  Every option has smaller birthday by Lemma~\ref{lem:birthday-options}, so
  the induction hypothesis identifies these entries literally with \(u\) and
  \(v\).  Thus the two option lists agree with those of \(a\), and hence
  \(a+0=a\) as game trees.  The proof of \(0+a=a\) is the same, with the
  summands interchanged.

  For commutativity, use the two-game induction scheme of
  Definition~\ref{def:game-inductions}.  By Definition~\ref{def:game-add}, a
  left option of \(a+b\) is either \(a^L+b\) or \(a+b^L\).  In the first
  case Lemma~\ref{lem:birthday-options} makes \((a^L,b)\) a smaller pair, and
  the induction hypothesis gives
  \[
    a^L+b\GameEq b+a^L;
  \]
  the right-hand side is a left option of \(b+a\).  In the second case the
  smaller pair \((a,b^L)\) gives
  \[
    a+b^L\GameEq b^L+a.
  \]
  The two right-option cases are identical, and interchanging \(a,b\) gives
  the converse matching.  Theorem~\ref{thm:option-extensionality} now yields
  \(a+b\GameEq b+a\).

  Associativity is stronger: it is literal equality of game trees.  Apply the
  one-game induction scheme of Definition~\ref{def:game-inductions}
  successively to \(a\), \(b\), and \(c\), generalizing the later variables at
  each stage.  The resulting induction hypotheses are
  \[
    \begin{array}{ll}
      ((a^L+b)+c)=a^L+(b+c) & (a^L\in A^L),\\
      ((a+b^L)+c)=a+(b^L+c) & (b^L\in B^L),\\
      ((a+b)+c^L)=a+(b+c^L) & (c^L\in C^L),
    \end{array}
  \]
  and the analogous three identities for right options.  Each recursive call
  is legitimate by Lemma~\ref{lem:birthday-options}; generalizing the later
  coordinates permits the other two games to remain unchanged.

  Expanding Definition~\ref{def:game-add} twice, the left option lists are
  \[
    \begin{aligned}
      ((a+b)+c)^L
       ={}&[\,(a^L+b)+c:a^L\in A^L\,]\\
          &{}\cup[\,(a+b^L)+c:b^L\in B^L\,]\\
          &{}\cup[\,(a+b)+c^L:c^L\in C^L\,],\\[2mm]
      (a+(b+c))^L
       ={}&[\,a^L+(b+c):a^L\in A^L\,]\\
          &{}\cup[\,a+(b^L+c):b^L\in B^L\,]\\
          &{}\cup[\,a+(b+c^L):c^L\in C^L\,].
    \end{aligned}
  \]
  We use \(\cup\) for the union of these finite option collections; formally, the underlying
  lists retain order and multiplicity.  On the left the three sublists are parenthesized as
  \((A^L\cup B^L)\cup C^L\), whereas on the right they are parenthesized as
  \(A^L\cup(B^L\cup C^L)\).  Associativity of finite-list concatenation first identifies these
  two parenthesizations.
  The lists then contain the same three families in the same order; only their
  entries differ.  The first, second, and third induction
  hypotheses above give pointwise equality on \(A^L,B^L,C^L\), respectively.
  Consequently the three corresponding sublists are equal, and therefore the
  complete left lists are equal.  Replacing \(A^L,B^L,C^L\) by
  \(A^R,B^R,C^R\) proves equality of the right lists in the same way.  The two
  games have identical left and right option lists, so
  \(((a+b)+c)=a+(b+c)\) literally.
\end{proof}
\begin{theorem}[Translation and monotonicity]
  \label{thm:game-add-order}
  \lean{Game.add_le_add_right,Game.add_right_cancel,Game.add_le_add,Game.add_equal,Game.add_lt_le,Game.add_le_lt,Game.add_lt_left_left,Game.add_lt_right_right}
  \leanok
  Addition preserves and reflects weak order in either variable.  In particular,
  \[
    a\GameLe b\quad\Longleftrightarrow\quad a+c\GameLe b+c,
  \]
  and
  \[
    a\GameLe c,\ b\GameLe d\Longrightarrow a+b\GameLe c+d.
  \]
  If either input comparison is strict and the other weak, the resulting comparison is strict.
  Addition also respects \(\GameEq\), i.e.
  \[
    a\GameEq b\quad\Longleftrightarrow\quad a+c\GameEq b+c.
  \]
\end{theorem}
\begin{proof}
  \leanok
  \uses{lem:birthday-options,def:game-order,def:game-equivalence,def:game-inductions,thm:game-order-calculus,def:game-add,thm:game-add-laws}
  We first prove simultaneously, in both directions, the equivalence
  \[
    a+c\GameLe b+c\quad\Longleftrightarrow\quad a\GameLe b.
  \]
  Use the three-game induction scheme of Definition~\ref{def:game-inductions}, with measure
  \(\bd(a)+\bd(b)+\bd(c)\), and use the comparison of
  Definition~\ref{def:game-order} throughout.

  Suppose first that \(a+c\GameLe b+c\).  If \(a^L\) is a left option of
  \(a\) and \(b\GameLe a^L\), the induction hypothesis for
  \((b,a^L,c)\) gives \(b+c\GameLe a^L+c\).  Transitivity from
  Theorem~\ref{thm:game-order-calculus}, together with the assumed comparison,
  gives \(a+c\GameLe a^L+c\).  This is impossible because
  Definition~\ref{def:game-add} makes \(a^L+c\) a left option of \(a+c\),
  while the left-option clause of the reflexive comparison
  \(a+c\GameLe a+c\) forbids such an inequality.  If \(b^R\) is a right
  option of \(b\) and \(b^R\GameLe a\), the induction hypothesis gives
  \(b^R+c\GameLe a+c\); composing with the assumed comparison contradicts the
  right-option clause of the reflexive comparison \(b+c\GameLe b+c\).  Thus
  \(a\GameLe b\).

  Conversely, assume \(a\GameLe b\).  By Definition~\ref{def:game-add}, the
  left options of \(a+c\) are \(a^L+c\) and \(a+c^L\), while the right
  options of \(b+c\) are \(b^R+c\) and \(b+c^R\).  If
  \(b+c\GameLe a^L+c\), the induction hypothesis for \((b,a^L,c)\) gives
  \(b\GameLe a^L\), contrary to the left-option clause of \(a\GameLe b\).
  If \(b+c\GameLe a+c^L\), the induction hypothesis for \((a,b,c^L)\)
  gives \(a+c^L\GameLe b+c^L\).  Transitivity would then give
  \(b+c\GameLe b+c^L\), contradicting the left-option clause of the reflexive
  comparison \(b+c\GameLe b+c\).  The right-option cases are dual.  A
  comparison \(b^R+c\GameLe a+c\) reflects to \(b^R\GameLe a\), which is
  forbidden by \(a\GameLe b\).  Finally, if
  \(b+c^R\GameLe a+c\), the induction hypothesis for \((a,b,c^R)\) gives
  \(a+c^R\GameLe b+c^R\); transitivity would imply
  \(a+c^R\GameLe a+c\), contrary to the right-option clause of the reflexive
  comparison \(a+c\GameLe a+c\).  This proves the reverse implication in the
  displayed equivalence.
  Every recursive call decreases because the coordinate replaced in it is an
  option, hence has smaller birthday by Lemma~\ref{lem:birthday-options}.

  Commutativity from Theorem~\ref{thm:game-add-laws} turns the displayed
  equivalence into the
  corresponding equivalence for addition on the left.  If
  \(a\GameLe c\) and \(b\GameLe d\), change first one summand and then the
  other, and compose the two inequalities using
  Theorem~\ref{thm:game-order-calculus}; this proves two-variable
  monotonicity.  Suppose, for example, that \(a\GameLt c\) and
  \(b\GameLe d\).  Weak monotonicity gives \(a+b\GameLe c+d\).  If the reverse
  comparison held, then
  \[
    c+b\GameLe c+d\GameLe a+b.
  \]
  Cancelling the common summand \(b\) by the displayed translation equivalence
  would give \(c\GameLe a\), contradicting \(a\GameLt c\).  Thus
  \(a+b\GameLt c+d\).  Strictness in the second variable is proved
  symmetrically.  Applying the weak result in both directions proves
  compatibility with \(\GameEq\).
\end{proof}

\begin{definition}[Negation]
  \label{def:game-neg}
  \lean{Game.neg,neg_left_def,neg_right_def,Game.neg_zero}
  \leanok
  \uses{def:game,def:birthday,lem:birthday-options}
  Negation interchanges the two option lists and recursively negates every option:
  \[
    -\{X^L\mid X^R\}=\{-X^R\mid -X^L\}.
  \]
  In particular, \(-0=0\) as a literal equality of game trees.
\end{definition}

\begin{theorem}[Negation reverses order]
  \label{thm:game-neg-order}
  \lean{Game.neg_le_neg,Game.neg_congr,Game.neg_lt_neg}
  \leanok
  For all games \(a,b\),
  \[
    a\GameLe b\Longleftrightarrow -b\GameLe-a,
    \qquad
    a\GameLt b\Longleftrightarrow -b\GameLt-a,
    \qquad
    a\GameEq b\Longleftrightarrow -b\GameEq-a.
  \]
\end{theorem}
\begin{proof}
  \leanok
  \uses{lem:birthday-options,def:game-order,def:game-equivalence,def:game-inductions,def:game-neg}
  Use the two-game induction scheme of Definition~\ref{def:game-inductions} to
  prove
  \[
    a\GameLe b\quad\Longleftrightarrow\quad -b\GameLe-a.
  \]
  Assume first \(a\GameLe b\).  By Definition~\ref{def:game-neg}, a left
  option of \(-b\) has the form \(-b^R\), and a right option of \(-a\) has
  the form \(-a^L\).  If \(-a\GameLe-b^R\), the induction hypothesis for
  \((b^R,a)\) gives \(b^R\GameLe a\), contradicting the right-option clause
  of \(a\GameLe b\) in Definition~\ref{def:game-order}.  Likewise, if
  \(-a^L\GameLe-b\), the induction hypothesis for \((b,a^L)\) gives
  \(b\GameLe a^L\), contradicting the left-option clause.  These are exactly
  the two clauses required for \(-b\GameLe-a\).

  Conversely, assume \(-b\GameLe-a\).  To verify \(a\GameLe b\), let
  \(a^L\) be a left option of \(a\).  A comparison \(b\GameLe a^L\) would,
  by the induction hypothesis for \((b,a^L)\), imply
  \(-a^L\GameLe-b\), contradicting the right-option clause of
  \(-b\GameLe-a\).  The argument for a right option \(b^R\) is the same:
  \(b^R\GameLe a\) would imply \(-a\GameLe-b^R\), contradicting the
  left-option clause.  Each induction step is valid by
  Lemma~\ref{lem:birthday-options}.

  Applying the displayed equivalence to both weak inequalities in
  Definition~\ref{def:game-equivalence} proves the stated biconditional for equivalence.
  Applying it to the forward weak
  inequality and to the excluded reverse inequality proves the strict-order
  equivalence.  Both conclusions use
  Definitions~\ref{def:game-order} and~\ref{def:game-equivalence}.
\end{proof}

\begin{theorem}[Additive inverses]
  \label{thm:game-add-inverse}
  \lean{Game.add_neg,Game.neg_add}
  \leanok
  Every game satisfies
  \[
    x+(-x)\GameEq0,
    \qquad
    (-x)+x\GameEq0.
  \]
\end{theorem}
\begin{proof}
  \leanok
  \uses{def:game,def:game-order,def:game-equivalence,def:game-inductions,lem:birthday-options,thm:game-order-calculus,def:game-add,thm:game-add-laws,def:game-neg}
  We argue by induction on \(\bd(x)\), using the one-game scheme of
  Definition~\ref{def:game-inductions}.  Put \(s=x+(-x)\).  We prove the two
  inequalities in the definition of \(s\GameEq0\) from
  Definition~\ref{def:game-equivalence}.  Since \(0=\{\,\mid\,\}\) by
  Definition~\ref{def:game}, the clauses involving options of \(0\) are
  vacuous.

  By Definitions~\ref{def:game-add} and~\ref{def:game-neg}, a left option of
  \(s\) is either
  \[
    \ell=x^L+(-x)
    \qquad\text{or}\qquad
    \ell=x+(-x^R).
  \]
  In the first case \(x^L+(-x^L)\) is a right option of \(\ell\); in the
  second case \(x^R+(-x^R)\) is a right option of \(\ell\).  By
  Lemma~\ref{lem:birthday-options} and the induction hypothesis, the indicated
  self-cancelling game is equivalent to \(0\), and in particular is weakly
  below \(0\).  Therefore \(0\GameLe\ell\) would violate the right-option
  clause of that comparison.  This verifies the left-option clause of
  \(s\GameLe0\).

  Similarly, a right option of \(s\) is either
  \[
    r=x^R+(-x)
    \qquad\text{or}\qquad
    r=x+(-x^L).
  \]
  The first has \(x^R+(-x^R)\) as a left option, and the second has
  \(x^L+(-x^L)\) as a left option.  The induction hypothesis makes that left
  option weakly above \(0\).  Thus \(r\GameLe0\) would violate its
  left-option clause.  This proves \(0\GameLe s\), and hence
  \(x+(-x)\GameEq0\).  Commutativity from
  Theorem~\ref{thm:game-add-laws} gives \((-x)+x\GameEq0\).
\end{proof}

\section{Surreal Number}

\begin{theorem}[Surreal games are closed under addition and negation]
  \label{thm:surreal-add-closure}
  \lean{Surreal.add_isSurreal,Surreal.neg_isSurreal,Surreal.add,Surreal.neg}
  \leanok
  If \(a,b\) are surreal representatives, then the games \(a+b\) and \(-a\) are surreal.
  Hence addition and negation define operations on surreal representatives.
\end{theorem}
\begin{proof}
  \leanok
  \uses{def:birthday,def:surreal-pair-induction,lem:birthday-options,def:game-order,def:is-surreal,thm:game-order-calculus,thm:between-options,def:surreal-representative,def:game-add,thm:game-add-order,def:game-neg,thm:game-neg-order}
  For addition, use the well-founded induction on pairs of surreal
  representatives from Definition~\ref{def:surreal-pair-induction}.  By
  Definition~\ref{def:game-add}, every option of \(a+b\) replaces exactly one
  of \(a,b\) by one of its options.  Lemma~\ref{lem:birthday-options} makes
  the resulting pair smaller, so the induction hypothesis proves that every
  option of \(a+b\) is surreal, as required by
  Definition~\ref{def:is-surreal}.

  It remains to separate a left option from a right option.  There are four
  cases:
  \[
    \begin{array}{c|c}
      \text{left option}&\text{right option}\\ \hline
      a^L+b&a^R+b\\
      a+b^L&a+b^R\\
      a^L+b&a+b^R\\
      a+b^L&a^R+b.
    \end{array}
  \]
  In the first case Theorem~\ref{thm:between-options} and strict transitivity
  from Theorem~\ref{thm:game-order-calculus} give
  \(a^L\GameLt a^R\), and strict translation invariance from
  Theorem~\ref{thm:game-add-order} gives
  \(a^L+b\GameLt a^R+b\).  The second case is identical with \(a,b\)
  interchanged.  In the third case,
  \(a^L\GameLt a\) and \(b\GameLe b^R\); in the fourth,
  \(a\GameLt a^R\) and \(b^L\GameLe b\).  The two mixed monotonicity
  statements of Theorem~\ref{thm:game-add-order} again put the chosen left
  option strictly below the chosen right option.  By
  Definition~\ref{def:game-order}, this strict inequality excludes the reverse
  weak comparison, which is exactly the separation condition in
  Definition~\ref{def:is-surreal}.

  For negation, induct on the birthday of the underlying game of the surreal
  representative \(a\).  Definition~\ref{def:game-neg} says that
  a left option of \(-a\) is \(-a^R\), while a right option is \(-a^L\).
  Theorem~\ref{thm:between-options} and strict transitivity from
  Theorem~\ref{thm:game-order-calculus} give
  \(a^L\GameLt a\GameLt a^R\), hence \(a^L\GameLt a^R\).  If
  \(-a^L\GameLe-a^R\), order reversal in
  Theorem~\ref{thm:game-neg-order} would give \(a^R\GameLe a^L\), contradicting
  that strict comparison.  Thus the left and right options of \(-a\) satisfy
  the separation condition of Definition~\ref{def:is-surreal}.
  Each negated option is surreal by the induction hypothesis and
  Lemma~\ref{lem:birthday-options}.  Thus \(-a\) satisfies
  Definition~\ref{def:is-surreal}.  Finally, Definition~\ref{def:surreal-representative}
  turns each game, together with the surrealness just proved, into the
  corresponding operation on representatives.
\end{proof}

\begin{definition}[Auxiliary quotient of all games]
  \label{def:game-quotient}
  \lean{Game.instSetoidGame,Game.GameQ,Game.q_eq}
  \leanok
  \uses{def:game-equivalence,thm:game-order-calculus}
  By Theorem~\ref{thm:game-order-calculus}, game equivalence is an equivalence relation.
  Define \(\Games/{\GameEq}\) to be the quotient of \(\Games\) by game equivalence, and
  write \([a]\) for the equivalence class of a game \(a\).  Thus \([a]=[b]\) if and only
  if \(a\GameEq b\).
\end{definition}

\begin{theorem}[Partially ordered additive commutative group]
  \label{thm:game-quotient-ordered-add-group}
  \lean{Game.q_zero,Game.q_add,Game.q_neg,Game.q_le,Game.q_lt,Game.instAddCommGroupGameQ,Game.instPartialOrderGameQ,Game.instIsOrderedAddMonoidGameQ}
  \leanok
  For finite games \(a,b\), define
  \[
    0=[0],\qquad [a]+[b]=[a+b],\qquad -[a]=[-a],
  \]
  and
  \[
    [a]\leq[b]\quad\Longleftrightarrow\quad a\GameLe b,
    \qquad
    [a]<[b]\quad\Longleftrightarrow\quad a\GameLt b.
  \]
  These formulas are independent of the chosen representatives.  With these operations and
  this order, \(\Games/{\GameEq}\) is a partially ordered additive commutative group, and
  addition is order preserving in each variable.
\end{theorem}
\begin{proof}
  \leanok
  \uses{def:game-quotient,def:game-order,def:game-equivalence,def:game-add,def:game-neg,thm:game-order-calculus,thm:game-add-laws,thm:game-add-order,thm:game-neg-order,thm:game-add-inverse}
  The addition and negation formulas are independent of the representatives by the
  congruence assertions in Theorems~\ref{thm:game-add-order}
  and~\ref{thm:game-neg-order}.

  The order is also independent of representatives.  Indeed, suppose
  \(a\GameEq a'\), \(b\GameEq b'\), and \(a\GameLe b\).  The two appropriate
  components of the equivalences give \(a'\GameLe a\) and
  \(b\GameLe b'\); transitivity from
  Theorem~\ref{thm:game-order-calculus} yields \(a'\GameLe b'\).  The converse
  follows symmetrically.  Reflexivity and transitivity therefore pass to
  classes.  If \([a]\leq[b]\) and \([b]\leq[a]\), then
  \(a\GameEq b\) by Definition~\ref{def:game-equivalence}, so \([a]=[b]\) by the
  definition of an equivalence class.  Thus the quotient order is a partial
  order.  The asserted characterization of strict order now follows from the
  definitions of \(\GameLt\) and strict order in a partial order.

  The zero, associativity, and commutativity laws descend from
  Theorem~\ref{thm:game-add-laws}, and the inverse law descends from
  Theorem~\ref{thm:game-add-inverse}; an equivalence of representatives is
  precisely equality of their classes.  Finally, translation monotonicity is
  the class-level form of Theorem~\ref{thm:game-add-order}.  Commutativity from
  Theorem~\ref{thm:game-add-laws} permits the common summand to be placed on
  either side.
\end{proof}

\begin{definition}[Short surreal numbers]
  \label{def:surreal-number}
  \lean{Surreal.SurrealNumber,Surreal.SurrealNumber.add,Surreal.SurrealNumber.neg}
  \leanok
  \uses{def:surreal-representative,thm:surreal-add-closure,thm:game-order-calculus,thm:game-add-order,thm:game-neg-order}
  The type \(\NoShort\) is the quotient of surreal representatives by \(a\GameEq b\).
  Addition, zero, one, and negation descend to this quotient.
\end{definition}

\begin{theorem}[Linear ordered additive commutative group]
  \label{thm:ordered-add-group}
  \lean{Surreal.instLinearOrderSurrealNumber,Surreal.instAddCommGroupSurrealNumber,Surreal.instIsOrderedAddMonoidSurrealNumber}
  \leanok
  The quotient \(\NoShort\) is a linearly ordered additive commutative group, and addition is
  compatible with its order.
\end{theorem}
\begin{proof}
  \leanok
  \uses{def:surreal-representative,thm:surreal-totality,thm:surreal-add-closure,def:game-quotient,thm:game-quotient-ordered-add-group,def:surreal-number}
  There is a natural injective map
  \[
    \iota:\NoShort\longrightarrow\Games/{\GameEq},
    \qquad \iota([s])=[s_{\mathrm{game}}].
  \]
  This is well defined and injective because both quotients identify representatives exactly
  when their underlying games are game equivalent.  Pull back the quotient order along
  \(\iota\), so \([s]\leq[t]\) precisely when
  \(s_{\mathrm{game}}\GameLe t_{\mathrm{game}}\).  Definition~\ref{def:surreal-number} and
  Theorem~\ref{thm:surreal-add-closure} show that the image contains zero and is closed under
  addition and negation.  Hence Theorem~\ref{thm:game-quotient-ordered-add-group} supplies the
  additive commutative group, partial order, and order compatibility on \(\NoShort\).
  Theorem~\ref{thm:surreal-totality} upgrades this partial order to a linear order.
\end{proof}


%%FILE blueprint/src/content.tex
% Shared content, used by both web.tex and print.tex.
\input{introduction}
\input{foundations}
\input{addition}
\input{multiplication}
\input{ordered_ring}


%%FILE blueprint/src/foundations.tex
\chapter{Games, order, and surreal representatives}

\section{Finite games and birthdays}

\begin{definition}[Short combinatorial game]
  \label{def:game}
  \lean{Game,Game.left,Game.right,Game.zero,Game.one}
  \leanok
  Let \(\Games\) denote the collection of all finite (or short) combinatorial games.  An
  element \(x\in\Games\) is a finite rooted tree written \(x=\{X^L\mid X^R\}\), where the
  left- and right-option collections \(X^L\) and \(X^R\) are finite lists of games.  The basic
  games are
  \[
    0=\{\,\mid\,\},\qquad 1=\{0\mid\,\}.
  \]
\end{definition}

\begin{definition}[Birthday rank]
  \label{def:birthday}
  \lean{Game.birthday}
  \leanok
  \uses{def:game}
  The rank \(\bd(x)\in\mathbb N\) is defined recursively by
  \[
    \bd\bigl(\{X^L\mid X^R\}\bigr)
      = 1+\max\bigl(\{\bd(u):u\in X^L\}\cup
                         \{\bd(v):v\in X^R\}\cup\{0\}\bigr).
  \]
  Thus our convention is \(\bd(0)=1\).
\end{definition}

\begin{lemma}[Options have smaller birthday]
  \label{lem:birthday-options}
  \lean{Game.birthday_lt_of_mem_options,Game.birthday_lt_left,Game.birthday_lt_right}
  \leanok
  If \(u\in X^L\) or \(u\in X^R\), then \(\bd(u)<\bd(x)\).
\end{lemma}
\begin{proof}
  \leanok
  \uses{def:birthday}
  Write \(x=\{X^L\mid X^R\}\), and let \(B\) be the finite collection of
  birthdays of all members of \(X^L\) and \(X^R\).  If \(u\) belongs to either
  option collection, then \(\bd(u)\in B\); in particular, \(B\) is nonempty
  and
  \[
    \bd(u)\leq \max B.
  \]
  By Definition~\ref{def:birthday}, \(\bd(x)=1+\max B\), where adjoining
  \(0\) makes the same formula valid when there are no options.  Hence
  \(\bd(u)<\bd(x)\).  The left- and right-option assertions are the two
  specializations of this argument.
\end{proof}

\section{The recursive order}

\begin{definition}[Game comparison]
  \label{def:game-order}
  \lean{Game.le,Game.lt}
  \leanok
  \uses{def:game,def:birthday,lem:birthday-options}
  The relation \(x\GameLe y\) is defined recursively by
  \[
    x\GameLe y
    \quad\Longleftrightarrow\quad
    \bigl(\forall x^L\in X^L,\ \neg(y\GameLe x^L)\bigr)
    \ \wedge\
    \bigl(\forall y^R\in Y^R,\ \neg(y^R\GameLe x)\bigr).
  \]
  Strict comparison is
  \[
    x\GameLt y \quad\Longleftrightarrow\quad
    x\GameLe y\ \wedge\ \neg(y\GameLe x).
  \]
\end{definition}

\begin{definition}[Game equivalence]
  \label{def:game-equivalence}
  \lean{Game.eq}
  \leanok
  \uses{def:game-order}
  Two games \(x,y\in\Games\) are \emph{equivalent}, written \(x\GameEq y\), when each is
  weakly below the other:
  \[
    x\GameEq y \quad\Longleftrightarrow\quad
    x\GameLe y\ \wedge\ y\GameLe x.
  \]
  This is the relation used later to identify games with the same game value.  Equivalent games
  need not be literally equal as rooted trees or have identical option lists.
\end{definition}

\begin{definition}[Well-founded induction schemes]
  \label{def:game-inductions}
  \lean{Game.R,Game.B,Game.T,Game.Q,Game.wf_R,Game.wf_B,Game.wf_T,Game.wf_Q}
  \leanok
  \uses{def:birthday}
  The development uses the following relations on one, two, three, and four games.  For
  \(q=(q_1,q_2,q_3,q_4)\) and \(q'=(q'_1,q'_2,q'_3,q'_4)\), set
  \[
    \begin{aligned}
      y\mathrel R x
        &\quad\Longleftrightarrow\quad \bd(y)<\bd(x),\\
      (a',b')\mathrel B(a,b)
        &\quad\Longleftrightarrow\quad
          \bd(a')+\bd(b')<\bd(a)+\bd(b),\\
      (a',b',c')\mathrel T(a,b,c)
        &\quad\Longleftrightarrow\quad
          \bd(a')+\bd(b')+\bd(c')<\bd(a)+\bd(b)+\bd(c),\\
      q'\mathrel Qq
        &\quad\Longleftrightarrow\quad
          \sum_{i=1}^{4}\bd(q'_i)<\sum_{i=1}^{4}\bd(q_i).
    \end{aligned}
  \]
  Each relation is well founded because it is the inverse image of the well-founded order
  \(<\) on \(\mathbb N\) under the corresponding sum-of-birthdays measure.  They will be used
  for well-founded induction on one, two, three, and four games, respectively; replacing any
  coordinate by an option strictly decreases the relevant measure.
\end{definition}

\begin{theorem}[Order calculus for games]
  \label{thm:game-order-calculus}
  \lean{Game.le_congr,Game.eq_congr,Game.le_trans',Game.le_trans,Game.eq_symm,Game.eq_trans,Game.lt_trans,Game.lt_of_lt_of_le,Game.lt_of_le_of_lt}
  \leanok
  For all games \(x,a,b,c\), reflexivity and transitivity are the assertions
  \[
    x\GameLe x,
    \qquad
    a\GameLe b\ \wedge\ b\GameLe c
      \quad\Longrightarrow\quad a\GameLe c.
  \]
  Consequently \(\GameEq\) is an equivalence relation, \(\GameLt\) is transitive, and strict
  and weak inequalities compose in either order.
\end{theorem}
\begin{proof}
  \leanok
  \uses{lem:birthday-options,def:game-order,def:game-equivalence,def:game-inductions}
  For reflexivity, use the one-game induction scheme of
  Definition~\ref{def:game-inductions}.  Suppose reflexivity is known for all
  games of birthday smaller than \(\bd(x)\), and unfold the comparison in
  Definition~\ref{def:game-order}.  If \(x^L\) is a left option and
  \(x\GameLe x^L\), then the left-option clause of this last comparison says
  \(\neg(x^L\GameLe x^L)\).  On the other hand,
  Lemma~\ref{lem:birthday-options} gives \(\bd(x^L)<\bd(x)\), so the induction
  hypothesis gives \(x^L\GameLe x^L\), a contradiction.  If \(x^R\) is a
  right option and \(x^R\GameLe x\), the right-option clause similarly says
  \(\neg(x^R\GameLe x^R)\), again contradicting the induction hypothesis.
  Thus \(x\GameLe x\).

  For transitivity, use the three-game induction scheme of
  Definition~\ref{def:game-inductions}.  Assume \(a\GameLe b\) and
  \(b\GameLe c\), and induct on
  \(\bd(a)+\bd(b)+\bd(c)\).  If \(a^L\) is a left option and
  \(c\GameLe a^L\), then the induction hypothesis applied to the triple
  \((b,c,a^L)\) gives \(b\GameLe a^L\).  This contradicts the left-option
  clause of \(a\GameLe b\) in Definition~\ref{def:game-order}.  Similarly, if
  \(c^R\) is a right option and \(c^R\GameLe a\), the induction hypothesis for
  \((c^R,a,b)\) gives \(c^R\GameLe b\), contradicting the right-option clause
  of \(b\GameLe c\).  Both recursive calls decrease because
  Lemma~\ref{lem:birthday-options} replaces one coordinate by an option of
  smaller birthday.  Hence \(a\GameLe c\).

  The remaining assertions follow directly from
  Definitions~\ref{def:game-order} and~\ref{def:game-equivalence}.  Symmetry and transitivity
  of \(\GameEq\)
  are obtained by composing the two weak inequalities in the appropriate
  directions.  For strict transitivity, compose the forward weak inequalities.
  If the reverse weak inequality held, transitivity would contradict one of
  the two strict hypotheses.  The mixed strict--weak laws follow from the same
  argument.
\end{proof}

\begin{theorem}[Equivalence from equivalent option lists]
  \label{thm:option-extensionality}
  \lean{Game.eq_of_equiv_options}
  \leanok
  Let \(a=\{A^L\mid A^R\}\) and \(b=\{B^L\mid B^R\}\) be games.  Suppose
  \[
    \begin{aligned}
      &\forall a^L\in A^L,\ \exists b^L\in B^L,\quad a^L\GameEq b^L,\\
      &\forall b^L\in B^L,\ \exists a^L\in A^L,\quad b^L\GameEq a^L,\\
      &\forall a^R\in A^R,\ \exists b^R\in B^R,\quad a^R\GameEq b^R,\\
      &\forall b^R\in B^R,\ \exists a^R\in A^R,\quad b^R\GameEq a^R.
    \end{aligned}
  \]
  Then \(a\GameEq b\).
\end{theorem}
\begin{proof}
  \leanok
  \uses{def:game-order,def:game-equivalence,thm:game-order-calculus}
  We first verify the two clauses of Definition~\ref{def:game-order} for
  \(a\GameLe b\).  Let \(a^L\) be a left option of \(a\), and suppose that
  \(b\GameLe a^L\).  Choose a left option \(b^L\) of \(b\) such that
  \(a^L\GameEq b^L\).  The forward inequality in this equivalence and
  transitivity from Theorem~\ref{thm:game-order-calculus} give
  \(b\GameLe b^L\).  This is impossible: reflexivity from the same theorem
  gives \(b\GameLe b\), whose left-option clause forbids
  \(b\GameLe b^L\).

  Next let \(b^R\) be a right option of \(b\), and suppose
  \(b^R\GameLe a\).  Choose a right option \(a^R\) of \(a\) such that
  \(b^R\GameEq a^R\).  The reverse inequality \(a^R\GameLe b^R\) in this
  equivalence and transitivity
  give \(a^R\GameLe a\), contrary to the right-option clause of the reflexive
  comparison \(a\GameLe a\).  Thus \(a\GameLe b\).  Interchanging \(a\) and
  \(b\) proves \(b\GameLe a\), and Definition~\ref{def:game-equivalence} then gives
  \(a\GameEq b\).
\end{proof}

\section{Surreal games}

\begin{definition}[Surreal game]
  \label{def:is-surreal}
  \lean{IsSurreal,IsSurreal.isSurreal_left,IsSurreal.isSurreal_right,IsSurreal.isSurreal_zero,IsSurreal.isSurreal_one}
  \leanok
  \uses{def:game,def:birthday,lem:birthday-options,def:game-order}
  A game \(x=\{X^L\mid X^R\}\) is surreal when every option is surreal and no right option is
  weakly below a left option:
  \[
    \forall x^L\in X^L\ \forall x^R\in X^R,
      \quad \neg(x^R\GameLe x^L).
  \]
  The games \(0\) and \(1\) are surreal, and every option of a surreal game is surreal.
\end{definition}

\begin{theorem}[A surreal game lies between its options]
  \label{thm:between-options}
  \lean{IsSurreal.left_lt,IsSurreal.lt_right,IsSurreal.xL_x_xR}
  \leanok
  If \(x\) is surreal, then
  \[
    \forall x^L\in X^L,\quad x^L\GameLt x,
    \qquad
    \forall x^R\in X^R,\quad x\GameLt x^R.
  \]
\end{theorem}
\begin{proof}
  \leanok
  \uses{lem:birthday-options,def:game-order,def:game-inductions,thm:game-order-calculus,def:is-surreal}
  We prove the assertion for left options; the right-option proof is its
  mirror.  Use the one-game induction scheme of
  Definition~\ref{def:game-inductions}.  Let \(x\) be surreal, and let
  \(x^L\) be a left option.  By Definition~\ref{def:is-surreal}, \(x^L\) is
  itself surreal.  We verify the two clauses of
  Definition~\ref{def:game-order} for \(x^L\GameLe x\).

  Let \(z\) be a left option of \(x^L\), and suppose \(x\GameLe z\).  Since
  \(\bd(x^L)<\bd(x)\) by Lemma~\ref{lem:birthday-options}, the induction
  hypothesis applied to the surreal game \(x^L\) gives \(z\GameLe x^L\).
  But the left-option clause of \(x\GameLe z\), applied to the option \(x^L\)
  of \(x\), forbids \(z\GameLe x^L\).  Next let \(x^R\) be a right option of
  \(x\).  The separation condition in Definition~\ref{def:is-surreal} is
  precisely \(\neg(x^R\GameLe x^L)\).  Hence \(x^L\GameLe x\).

  It remains to exclude the reverse comparison.  If \(x\GameLe x^L\), its
  left-option clause gives \(\neg(x^L\GameLe x^L)\), contradicting
  reflexivity from Theorem~\ref{thm:game-order-calculus}.  Thus
  \(x^L\GameLt x\) by Definition~\ref{def:game-order}.  The mirrored argument
  proves \(x\GameLt x^R\) for every right option \(x^R\).
\end{proof}

\begin{theorem}[Totality and trichotomy for surreal games]
  \label{thm:surreal-totality}
  \lean{IsSurreal.le_of_not_le,IsSurreal.totality,IsSurreal.trichotomy}
  \leanok
  If \(x\) and \(y\) are surreal, then
  \[
    x\GameLe y\ \vee\ y\GameLe x,
  \]
  and one of the three alternatives \(x\GameLt y\), \(x\GameEq y\), or
  \(y\GameLt x\) holds.
\end{theorem}
\begin{proof}
  \leanok
  \uses{def:game-order,def:game-equivalence,def:is-surreal,thm:game-order-calculus,thm:between-options}
  Suppose \(\neg(x\GameLe y)\).  Negating the two universal clauses in
  Definition~\ref{def:game-order} produces one of two witnesses: either there
  is a left option \(x^L\) such that \(y\GameLe x^L\), or there is a right
  option \(y^R\) such that \(y^R\GameLe x\).  In the first case,
  Theorem~\ref{thm:between-options} gives \(x^L\GameLe x\); in the second it
  gives \(y\GameLe y^R\).  Transitivity from
  Theorem~\ref{thm:game-order-calculus} therefore yields \(y\GameLe x\) in
  either case.  Hence
  \[
    \neg(x\GameLe y)\quad\Longrightarrow\quad y\GameLe x,
  \]
  which proves totality.

  For trichotomy, choose a direction supplied by totality and then distinguish
  whether the reverse weak comparison also holds.  If both hold, then
  \(x\GameEq y\); if exactly one holds, then the corresponding comparison is
  strict.  These are precisely the definitions in
  Definitions~\ref{def:game-order} and~\ref{def:game-equivalence}.
\end{proof}

\begin{definition}[Surreal representatives]
  \label{def:surreal-representative}
  \lean{Surreal,Surreal.eq,Surreal.sr_zero,Surreal.sr_one,Surreal.leftOption,Surreal.rightOption}
  \leanok
  \uses{def:game,def:game-equivalence,def:is-surreal}
  Let \(\SurrealReps\) denote the collection of all surreal representatives:
  \[
    \SurrealReps
      =\bigl\{\,s=(s_{\mathrm{game}},s_{\mathrm{proof}})
        \bigm|
        s_{\mathrm{game}}\in\Games,\quad
        s_{\mathrm{proof}}\text{ is a proof of }
        \operatorname{IsSurreal}(s_{\mathrm{game}})\,\bigr\}.
  \]
  For \(s,s'\in\SurrealReps\), representative equivalence is induced by game equivalence:
  \[
    s\GameEq s'
      \quad\Longleftrightarrow\quad
    s_{\mathrm{game}}\GameEq s'_{\mathrm{game}}.
  \]
  The surreal games \(0\) and \(1\) give canonical elements of \(\SurrealReps\).  If
  \(s\in\SurrealReps\), then each left or right option of \(s_{\mathrm{game}}\), equipped with
  the inherited surrealness proof from Definition~\ref{def:is-surreal}, is again an element of
  \(\SurrealReps\).  Equivalent representatives are related by \(\GameEq\), but are not yet
  identified as equal elements of a quotient.
\end{definition}

\begin{definition}[Pair induction for surreal representatives]
  \label{def:surreal-pair-induction}
  \lean{Surreal.BiSurreal,Surreal.U,Surreal.wf_U,Surreal.U_left₁,Surreal.U_left₂,Surreal.U_right₁,Surreal.U_right₂}
  \leanok
  \uses{def:birthday,lem:birthday-options,def:surreal-representative}
  Let \(a,b\in\SurrealReps\), and write \(\{a,b\}\) for the formal two-coordinate object with
  entries \(a\) and \(b\).  For \(a',b'\in\SurrealReps\), define
  \[
    \{a',b'\}\mathrel{\vartriangleleft}\{a,b\}
      \quad\Longleftrightarrow\quad
    \bd(a'_{\mathrm{game}})+\bd(b'_{\mathrm{game}})
      <\bd(a_{\mathrm{game}})+\bd(b_{\mathrm{game}}).
  \]
  This relation is well founded.  Replacing either coordinate by any left or right option
  produces a smaller pair, by Lemma~\ref{lem:birthday-options}.
\end{definition}


%%FILE blueprint/src/introduction.tex
\chapter{Overview}

This blueprint records the mathematical dependency structure of the
\href{https://github.com/kdwong/surreal}{surreal Lean repository}.  It follows the declarations
and proofs in that development rather than presenting a separate textbook construction.

Every game considered here has finitely many options at each node and finite birthday.
Accordingly, the quotient constructed here is the ring of
\emph{short} surreal numbers, denoted in this document by \(\NoShort\).  This is the dyadic
part of Conway's full class of surreal numbers.  No claim about proper-class-sized or
transfinite birthdays is made.

The principal dependency path is
\[
\begin{aligned}
  &\text{finite games and birthdays}
  \longrightarrow \text{recursive game order}
  \longrightarrow \text{surreal totality} \\
  &\quad\longrightarrow \text{ordered additive quotient},\\[2mm]
  &\text{recursive product}
  \longrightarrow \text{Conway A/B/C induction}
  \longrightarrow \text{product closure and congruence}\\
  &\quad\longrightarrow \text{commutative ring}
  \longrightarrow \text{positive-product theorem}
  \longrightarrow \text{linear ordered commutative ring}.
\end{aligned}
\]

On the web version, each formalized item links to its Lean declaration.  The dependency
graph uses dashed edges for dependencies occurring in a statement or construction and solid
edges for proof dependencies.  The graph retains every declared edge, including transitively
redundant ones, so the complete audited dependency path remains visible.

All nodes marked as formalized were checked against the current local formal development.  A
green node means that the corresponding theorem has been verified; the prose proof explains
its mathematical argument and its place in the dependency path.


%%FILE blueprint/src/multiplication.tex
\chapter{Multiplication and the Conway induction}

\section{The recursive product}

\begin{definition}[Product and product-option expression]
  \label{def:game-mul}
  \lean{Game.mul,Game.mulOpt4,mem_mul_left,mem_mul_right}
  \leanok
  \uses{def:game,def:birthday,lem:birthday-options,def:game-inductions,def:game-add,def:game-neg}
  Put
  \[
    \mulopt(x',y;x,y')=x'y+xy'-x'y'.
  \]
  For \(x=\{X^L\mid X^R\}\) and \(y=\{Y^L\mid Y^R\}\), multiplication is the recursive cut
  \[
    xy=
    \left\{
      \begin{array}{ll}
        \mulopt(x^L,y;x,y^L),&\mulopt(x^R,y;x,y^R)
      \end{array}
      \ \middle|\
      \begin{array}{ll}
        \mulopt(x^L,y;x,y^R),&\mulopt(x^R,y;x,y^L)
      \end{array}
    \right\},
  \]
  with all choices of the displayed options.  By
  Lemma~\ref{lem:birthday-options}, each selected option has smaller birthday
  than its parent.  Hence the two-game induction relation of
  Definition~\ref{def:game-inductions} decreases in every recursive call.
\end{definition}

\begin{theorem}[Zero, one, and commutativity for game multiplication]
  \label{thm:game-mul-laws}
  \lean{Game.mul_zero_eq,Game.zero_mul_eq,Game.mul_zero,Game.zero_mul,Game.mul_comm,Game.mul_one_eq,Game.mul_one,Game.one_mul}
  \leanok
  For all games \(a,b\),
  \[
    a0\GameEq0,\qquad 0a\GameEq0,\qquad ab\GameEq ba,
    \qquad a1\GameEq a,\qquad 1a\GameEq a.
  \]
  In fact, multiplication by zero on either side and multiplication on the right by one
  satisfy literal game-tree equalities.
\end{theorem}
\begin{proof}
  \leanok
  \uses{def:game,lem:birthday-options,def:game-inductions,thm:game-order-calculus,thm:option-extensionality,thm:game-add-laws,thm:game-add-order,def:game-neg,thm:game-neg-order,def:game-mul}
  In \(a0\), the second factor has no options.  Expanding the recursive cut in
  Definition~\ref{def:game-mul} therefore gives two empty option lists, so
  \(a0=0\) as literal game trees.
  The same direct expansion of \(0a\) has no option in any of the four
  families, because the first factor has no options; hence \(0a=0\) literally
  as well.

  For commutativity, use the two-game well-founded induction of
  Definition~\ref{def:game-inductions}; every recursive pair is smaller by
  Lemma~\ref{lem:birthday-options}.  Split an option according to the four
  families in Definition~\ref{def:game-mul}.  For a same-side example, the
  \(LL\) option of
  \(ab\) determined by \((a^L,b^L)\) is
  \[
    \mulopt(a^L,b;a,b^L)=a^Lb+ab^L-a^Lb^L.
  \]
  The \(LL\) option of \(ba\) determined by \((b^L,a^L)\) is
  \[
    \mulopt(b^L,a;b,a^L)=b^La+ba^L-b^La^L.
  \]
  The induction hypotheses for the three smaller pairs
  \((a^L,b)\), \((a,b^L)\), and \((a^L,b^L)\) identify, respectively,
  \(a^Lb\GameEq ba^L\), \(ab^L\GameEq b^La\), and
  \(a^Lb^L\GameEq b^La^L\).  Additive commutativity is
  Theorem~\ref{thm:game-add-laws}, compatibility of addition with
  \(\GameEq\) is part of Theorem~\ref{thm:game-add-order}, and compatibility
  of negation with \(\GameEq\) is part of
  Theorem~\ref{thm:game-neg-order}.  These three results therefore make the
  two displayed options equivalent.  The \(RR\) case is the identical
  calculation with
  \(a^L,b^L\) replaced by \(a^R,b^R\):
  \[
    \mulopt(a^R,b;a,b^R)\GameEq\mulopt(b^R,a;b,a^R).
  \]

  A mixed-side option changes family when the factors are swapped.  The
  \(LR\) option
  \[
    \mulopt(a^L,b;a,b^R)=a^Lb+ab^R-a^Lb^R
  \]
  is matched not with another \(LR\) option, but with the \(RL\) option
  \[
    \mulopt(b^R,a;b,a^L)=b^Ra+ba^L-b^Ra^L
  \]
  of \(ba\).  The induction hypotheses for \((a^L,b)\), \((a,b^R)\), and
  \((a^L,b^R)\) identify its three terms after the same additive
  rearrangement.  Dually,
  \[
    \mulopt(a^R,b;a,b^L)\GameEq\mulopt(b^L,a;b,a^R),
  \]
  so an \(RL\) option of \(ab\) matches an \(LR\) option of \(ba\).
  Repeating these matches from \(ba\) back to \(ab\), for both the left and
  right option lists, supplies the four hypotheses of
  Theorem~\ref{thm:option-extensionality}; hence \(ab\GameEq ba\).  The two
  literal zero identities already prove both displayed zero laws.

  Finally, prove the literal equality \(a1=a\) by the one-game induction of
  Definition~\ref{def:game-inductions}.  By Definition~\ref{def:game},
  \(1=\{0\mid\}\).  Thus an \(LL\) option of \(a1\) is
  \[
    \mulopt(a^L,1;a,0)=a^L1+a0-a^L0,
  \]
  and an \(RL\) option is the same expression with \(a^R\) in place of
  \(a^L\); the other two families are empty.  The induction hypothesis gives
  \(a^L1=a^L\) and \(a^R1=a^R\), while the zero-product identity established
  above, the additive zero laws in Theorem~\ref{thm:game-add-laws}, and the
  identity \(-0=0\) in Definition~\ref{def:game-neg} reduce the displayed
  expressions to \(a^L\) and \(a^R\).  Hence the left and right option lists of \(a1\) are
  literally those of \(a\).  The commutativity conclusion already proved then
  gives \(1a\GameEq a\).
\end{proof}
\begin{theorem}[Distributivity on games]
  \label{thm:game-distrib}
  \lean{Game.mul_distrib_tri,Game.mul_distrib,Game.mul_distrib_right}
  \leanok
  For all games \(a,b,c\),
  \[
    a(b+c)\GameEq ab+ac,
    \qquad
    (a+b)c\GameEq ac+bc.
  \]
\end{theorem}
\begin{proof}
  \leanok
  \uses{def:game-inductions,lem:birthday-options,thm:game-order-calculus,thm:option-extensionality,def:game-add,thm:game-add-order,thm:game-neg-order,def:game-quotient,thm:game-quotient-ordered-add-group,def:game-mul,thm:game-mul-laws}
  Prove the first equivalence by the three-game well-founded induction of
  Definition~\ref{def:game-inductions}.  Lemma~\ref{lem:birthday-options}
  shows that every recursive triple is smaller.  For an option of
  \(a(b+c)\), Definition~\ref{def:game-mul} first distinguishes whether the
  changed coordinate is \(a\) or \(b+c\); in the second case,
  Definition~\ref{def:game-add} distinguishes whether the changed coordinate
  is \(b\) or \(c\).  The matching option of \(ab+ac\) changes the same
  coordinate in the corresponding summand.  Conversely, an option of the sum
  changes either \(ab\) or \(ac\), and Definition~\ref{def:game-mul}
  determines the corresponding option on the left.  This gives four
  left/right matching assertions, two in each direction.

  For example, let \(a^L\) and \(b^L\) be left options.  The corresponding
  \(LL\) option on the left is
  \[
    L=\mulopt(a^L,b+c;a,b^L+c).
  \]
  The matching left option on the right is
  \[
    L'=\mulopt(a^L,b;a,b^L)+ac.
  \]
  The three smaller induction hypotheses give
  \[
    \begin{aligned}
      a^L(b+c)&\GameEq a^Lb+a^Lc,\\
      a(b^L+c)&\GameEq ab^L+ac,\\
      a^L(b^L+c)&\GameEq a^Lb^L+a^Lc.
    \end{aligned}
  \]
  Substitution in the definition of \(\mulopt\) gives
  \[
    L\GameEq
      (a^Lb+a^Lc)+(ab^L+ac)-(a^Lb^L+a^Lc)
      \GameEq \mulopt(a^L,b;a,b^L)+ac=L'.
  \]
  The cancellation in the second step takes place in the auxiliary quotient of
  Definition~\ref{def:game-quotient}, with its ordered additive group structure from
  Theorem~\ref{thm:game-quotient-ordered-add-group}.  Every remaining
  membership branch uses one of the same two algebraic templates, according
  as the changed summand is \(b\) or \(c\); reverse-direction branches use
  symmetry of the resulting equivalence.  Left and right options are placed
  as prescribed by Definition~\ref{def:game-mul}.  Compatibility with sums and negatives
  and additive-group normalization follow from
  Theorem~\ref{thm:game-quotient-ordered-add-group}, which puts each pair of signed sums
  in the same form.  By the
  quotient relation in
  Definition~\ref{def:game-quotient}, equality of their classes is precisely
  game equivalence.  The four matches now satisfy
  Theorem~\ref{thm:option-extensionality}, proving
  \(a(b+c)\GameEq ab+ac\).

  Finally,
  \[
    (a+b)c\GameEq c(a+b)\GameEq ca+cb\GameEq ac+bc.
  \]
  The first equivalence uses product commutativity from
  Theorem~\ref{thm:game-mul-laws}; the middle equivalence is the
  left-distributive result just proved; and the last uses product
  commutativity inside the two summands together with addition congruence from
  Theorem~\ref{thm:game-add-order}.
\end{proof}

\begin{lemma}[Multiplication and negation]
  \label{lem:game-mul-neg}
  \lean{Game.mul_neg,Game.neg_mul}
  \leanok
  For all games \(u,v\),
  \[
    u(-v)\GameEq-(uv),
    \qquad
    (-u)v\GameEq-(uv).
  \]
\end{lemma}
\begin{proof}
  \leanok
  \uses{def:game-inductions,lem:birthday-options,def:game-neg,def:game-mul,
    thm:game-order-calculus,thm:option-extensionality,thm:game-neg-order,
    def:game-quotient,thm:game-quotient-ordered-add-group,thm:game-mul-laws}
  Prove the first equivalence by the two-game induction of
  Definition~\ref{def:game-inductions}.  Let \(u'\) be either a left or a
  right option of \(u\), and let \(v'\) be either a left or a right option of
  \(v\).  The three smaller induction hypotheses are
  \[
    u'(-v)\GameEq-(u'v),\qquad
    u(-v')\GameEq-(uv'),\qquad
    u'(-v')\GameEq-(u'v').
  \]
  In the auxiliary quotient, these hypotheses give
  \[
    \begin{aligned}
      [\mulopt(u',-v;u,-v')]
        &=[u'(-v)]+[u(-v')]-[u'(-v')]\\
        &=-[u'v]-[uv']+[u'v']\\
        &=-\bigl([u'v]+[uv']-[u'v']\bigr)\\
        &=[-\mulopt(u',v;u,v')].
    \end{aligned}
  \]
  Equality of the first and last classes means
  \[
    \mulopt(u',-v;u,-v')
      \GameEq-\mulopt(u',v;u,v').
    \tag{*}
  \]

  Definition~\ref{def:game-neg} gives
  \(({-v})^L=-v^R\) and \(({-v})^R=-v^L\).  Consequently, applying \((*)\)
  to the indicated choices gives the two left-option matches
  \[
    \begin{aligned}
      \mulopt(u^L,-v;u,-v^R)
        &\GameEq-\mulopt(u^L,v;u,v^R),\\
      \mulopt(u^R,-v;u,-v^L)
        &\GameEq-\mulopt(u^R,v;u,v^L),
    \end{aligned}
  \]
  and the two right-option matches
  \[
    \begin{aligned}
      \mulopt(u^L,-v;u,-v^L)
        &\GameEq-\mulopt(u^L,v;u,v^L),\\
      \mulopt(u^R,-v;u,-v^R)
        &\GameEq-\mulopt(u^R,v;u,v^R).
    \end{aligned}
  \]
  The expressions on the right in the first display are precisely the
  negatives of the two right options of \(uv\), hence are left options of
  \(-(uv)\); those in the second display are the negatives of the two left
  options of \(uv\), hence are right options of \(-(uv)\).  Reading the same
  matches in reverse supplies the converse directions.  Thus
  Theorem~\ref{thm:option-extensionality} proves
  \(u(-v)\GameEq-(uv)\).

  Finally,
  \[
    (-u)v
      \GameEq v(-u)
      \GameEq-(vu)
      \GameEq-(uv).
  \]
  The first and last steps use product commutativity from
  Theorem~\ref{thm:game-mul-laws}, with negation congruence from
  Theorem~\ref{thm:game-neg-order} in the last step; the middle step is the
  first equivalence.
\end{proof}

\section{A simultaneous well-founded proof}

\begin{definition}[Dershowitz--Manna goal order]
  \label{def:dm-goals}
  \lean{Goal,μA,μB,μGoal,GoalLT,wfGoal}
  \leanok
  \uses{def:birthday,lem:birthday-options}
  Three kinds of induction goal are used:
  \[
    A(x,y),\qquad B(x_1,x_2,y),\qquad C(x_1,x_2,y).
  \]
  Their measures are the multisets
  \[
    \mu_A(x,y)=\{\bd(x),\bd(y)\},\qquad
    \mu_B(x_1,x_2,y)=\mu_C(x_1,x_2,y)
      =\{\bd(x_1),\bd(x_2),\bd(y)\}.
  \]
  The goal order is the pullback of the Dershowitz--Manna multiset extension
  \(<_{\rm DM}\) of \(<\) on \(\mathbb N\) along these measures.  For example,
  if \(a'<a\) and \(a_1,a_2<a\), then
  \[
    \{a',b\}<_{\rm DM}\{a,b\},\qquad
    \{a',b,c\}<_{\rm DM}\{a,b,c\},
  \]
  \[
    \{a_1,a_2,b\}<_{\rm DM}\{a,b\},\qquad
    \{a,b\}<_{\rm DM}\{a,b,c\}.
  \]
  The first two comparisons replace one birthday by a smaller birthday, the
  third replaces one birthday by two smaller birthdays, and the fourth deletes
  one birthday.  Lemma~\ref{lem:birthday-options}, transitivity, and permutation
  invariance of multisets reduce the recursive dependencies in the
  \(A\)-, \(B\)-, and \(C\)-cases to these examples.  Since \(<\) on
  \(\mathbb N\) is well founded, so is the resulting goal order.
\end{definition}

\begin{definition}[Cross inequalities]
  \label{def:cross-conditions}
  \lean{CLeft,CRight,MulCrossLe,MulCrossLt}
  \leanok
  \uses{def:game-order,def:game-add,def:game-mul}
  For \(x_1\GameLt x_2\), the two forms needed to compare product options are
  \[
    \begin{aligned}
      C_L(x_1,x_2;y,y^L)&:\quad
        x_1y+x_2y^L\GameLt x_1y^L+x_2y,\\
      C_R(x_1,x_2;y,y^R)&:\quad
        x_1y^R+x_2y\GameLt x_1y+x_2y^R.
    \end{aligned}
  \]
  For arbitrary games \(x_1,x_2,y_1,y_2\), define the weak and strict
  four-product relations by
  \[
    \begin{aligned}
      \operatorname{Cross}_{\GameLe}(x_1,x_2;y_1,y_2)&:\quad
        x_1y_2+x_2y_1\GameLe x_1y_1+x_2y_2,\\
      \operatorname{Cross}_{\GameLt}(x_1,x_2;y_1,y_2)&:\quad
        x_1y_2+x_2y_1\GameLt x_1y_1+x_2y_2.
    \end{aligned}
  \]
  Thus \(C_L(x_1,x_2;y,y^L)\) is
  \(\operatorname{Cross}_{\GameLt}(x_1,x_2;y^L,y)\), while
  \(C_R(x_1,x_2;y,y^R)\) is
  \(\operatorname{Cross}_{\GameLt}(x_1,x_2;y,y^R)\).
\end{definition}

\begin{theorem}[Cross-inequality calculus]
  \label{thm:cross-calculus}
  \lean{MulCrossLt.swap,MulCrossLt.trans,MulCrossLt.congr_right,MulCrossLt.congr_left,mulOpt4_xslot_mono_le_of_cross,mulOpt4_xslot_antitone_le_of_cross,mulOpt4_yslot_mono_le_of_cross,mulOpt4_yslot_antitone_le_of_cross,mulOpt4_xslot_mono_lt_of_cross,mulOpt4_xslot_antitone_lt_of_cross,mulOpt4_yslot_mono_lt_of_cross,mulOpt4_yslot_antitone_lt_of_cross}
  \leanok
  Let \(\mathrel{\bowtie}\) denote either \(\GameLe\) or \(\GameLt\).  If
  \(\operatorname{Cross}_{\bowtie}(a,c;d,e)\) holds, then
  \[
    \begin{array}{ll}
      \text{(a)}&
        \mulopt(a,e;b,d)\mathrel{\bowtie}\mulopt(c,e;b,d),\\
      \text{(b)}&
        \mulopt(c,d;b,e)\mathrel{\bowtie}\mulopt(a,d;b,e),\\
      \text{(c)}&
        \mulopt(a,b;c,d)\mathrel{\bowtie}\mulopt(a,b;c,e),\\
      \text{(d)}&
        \mulopt(c,b;a,e)\mathrel{\bowtie}\mulopt(c,b;a,d).
    \end{array}
  \]
  For the strict relation, the transpose, composition, and two endpoint-replacement rules are
  \[
    \begin{aligned}
      \text{(e)}\quad&
      \operatorname{Cross}_{\GameLt}(a,c;d,e)
        \Longrightarrow \operatorname{Cross}_{\GameLt}(d,e;a,c),\\
      \text{(f)}\quad&
      \operatorname{Cross}_{\GameLt}(a,b;d,e)\ \wedge
      \operatorname{Cross}_{\GameLt}(b,c;d,e)
        \Longrightarrow \operatorname{Cross}_{\GameLt}(a,c;d,e),\\
      \text{(g)}\quad&
      \operatorname{Cross}_{\GameLt}(a,b;d,e)\ \wedge
      bd\GameEq cd\ \wedge\ be\GameEq ce
        \Longrightarrow \operatorname{Cross}_{\GameLt}(a,c;d,e),\\
      \text{(h)}\quad&
      ad\GameEq bd\ \wedge\ ae\GameEq be\ \wedge
      \operatorname{Cross}_{\GameLt}(b,c;d,e)
        \Longrightarrow \operatorname{Cross}_{\GameLt}(a,c;d,e).
    \end{aligned}
  \]
\end{theorem}
\begin{proof}
  \leanok
  \uses{def:game-add,thm:game-add-laws,def:game-mul,thm:game-mul-laws,def:game-quotient,thm:game-quotient-ordered-add-group,def:cross-conditions}
  Work in the quotient \(\mathcal G/{\GameEq}\) of
  Definition~\ref{def:game-quotient}, equipped with the ordered additive group structure of
  Theorem~\ref{thm:game-quotient-ordered-add-group}.  The hypothesis is
  \[
    [ae]+[cd]\mathrel{\bowtie}[ad]+[ce].
  \]
  Adding the same class to both sides and rearranging the additive group gives,
  respectively,
  \[
    \begin{aligned}
      {}[ae]+[bd]-[ad]&\mathrel{\bowtie}[ce]+[bd]-[cd],\\
      {}[cd]+[be]-[ce]&\mathrel{\bowtie}[ad]+[be]-[ae],\\
      {}[ab]+[cd]-[ad]&\mathrel{\bowtie}[ab]+[ce]-[ae],\\
      {}[cb]+[ae]-[ce]&\mathrel{\bowtie}[cb]+[ad]-[cd].
    \end{aligned}
  \]
  For example, Definition~\ref{def:game-mul} gives
  \[
    \mulopt(a,e;b,d)=ae+bd-ad,
    \qquad
    \mulopt(c,e;b,d)=ce+bd-cd.
  \]
  Hence the first rearranged inequality is exactly
  \[
    [\mulopt(a,e;b,d)]\mathrel{\bowtie}[\mulopt(c,e;b,d)],
  \]
  which is the class-level form of rule \textup{(a)}.
  The other three lines give rules \textup{(b)}--\textup{(d)} in the same way.
  The argument works for both weak and strict order because translation preserves both in
  the ordered additive group.

  We prove rule \textup{(g)},
  \[
    \operatorname{Cross}_{\GameLt}(a,b;d,e)
      \ \wedge\ bd\GameEq cd
      \ \wedge\ be\GameEq ce
      \quad\Longrightarrow\quad
    \operatorname{Cross}_{\GameLt}(a,c;d,e);
  \]
  the other strict rules \textup{(e)}, \textup{(f)}, and \textup{(h)} are analogous.
  The first hypothesis means
  \[
    ae+bd\GameLt ad+be.
  \]
  By Theorem~\ref{thm:game-quotient-ordered-add-group}, this is the strict
  quotient inequality
  \[
    [ae]+[bd]<[ad]+[be].
  \]
  The two equivalence hypotheses give the class equalities
  \[
    [bd]=[cd],
    \qquad
    [be]=[ce].
  \]
  Substituting them into the strict inequality gives
  \[
    [ae]+[cd]<[ad]+[ce].
  \]
  Reflecting strict order from the quotient yields
  \[
    ae+cd\GameLt ad+ce,
  \]
  which is exactly
  \(\operatorname{Cross}_{\GameLt}(a,c;d,e)\).
\end{proof}

\begin{lemma}[Terminal product-option bounds]
  \label{lem:mulopt-bounds}
  \lean{mulOpt4_LL_lt_product,mulOpt4_RR_lt_product,product_lt_mulOpt4_LR,product_lt_mulOpt4_RL}
  \leanok
  The following statements hold:
  \[
    \begin{array}{ll}
      (LL)&ay^L\GameEq by^L,\ C_L(a^L,b;y,y^L)
        \Longrightarrow \mulopt(a^L,y;a,y^L)\GameLt by,\\[1mm]
      (RR)&ay^R\GameEq by^R,\ C_R(b,a^R;y,y^R)
        \Longrightarrow \mulopt(a^R,y;a,y^R)\GameLt by,\\[1mm]
      (LR)&by^R\GameEq ay^R,\ C_R(b^L,a;y,y^R)
        \Longrightarrow ay\GameLt\mulopt(b^L,y;b,y^R),\\[1mm]
      (RL)&by^L\GameEq ay^L,\ C_L(a,b^R;y,y^L)
        \Longrightarrow ay\GameLt\mulopt(b^R,y;b,y^L).
    \end{array}
  \]
\end{lemma}
\begin{proof}
  \leanok
  \uses{def:game-mul,def:game-quotient,thm:game-quotient-ordered-add-group,def:cross-conditions}
  We prove only the \(LL\) statement.  Its hypotheses are
  \[
    ay^L\GameEq by^L
    \quad\text{and}\quad
    C_L(a^L,b;y,y^L):
    a^Ly+by^L\GameLt a^Ly^L+by.
  \]
  In the quotient \(\mathcal G/{\GameEq}\), add \(-[a^Ly^L]\) to both sides and use
  \([by^L]=[ay^L]\).  This gives
  \[
    [a^Ly]+[ay^L]-[a^Ly^L]<[by],
  \]
  whose left side is \([\mulopt(a^L,y;a,y^L)]\) by
  Definition~\ref{def:game-mul}.  The quotient order in
  Theorem~\ref{thm:game-quotient-ordered-add-group} therefore gives
  \[
    \mulopt(a^L,y;a,y^L)\GameLt by.
  \]
  The \(RR\), \(LR\), and \(RL\) statements follow from the same quotient
  rearrangement with their displayed equivalence and cross-inequality hypotheses.
\end{proof}

\begin{theorem}[Simultaneous Conway theorem]
  \label{thm:abc-main}
  \lean{Conway.Holds,Conway.ABC_main}
  \leanok
  The following assertions hold simultaneously for all games satisfying the
  displayed hypotheses:
  \begin{enumerate}
    \item \(A(x,y)\): if \(x,y\) are surreal, then \(xy\) is surreal;
    \item \(B(x_1,x_2,y)\): if \(x_1\GameEq x_2\) and all three games are surreal, then
      \(x_1y\GameEq x_2y\);
    \item \(C(x_1,x_2,y)\): if \(x_1\GameLt x_2\) and all three games are surreal, then
      \(C_L(x_1,x_2;y,y^L)\) holds for every \(y^L\), and
      \(C_R(x_1,x_2;y,y^R)\) holds for every \(y^R\).
  \end{enumerate}
\end{theorem}
\begin{proof}
  \leanok
  \uses{def:game-order,def:game-equivalence,def:is-surreal,def:game-mul,def:game-quotient,thm:game-quotient-ordered-add-group,def:dm-goals,def:cross-conditions,thm:game-order-calculus,thm:between-options,thm:surreal-totality,thm:surreal-add-closure,thm:game-mul-laws,thm:cross-calculus,lem:mulopt-bounds}
  Apply well-founded induction with the goal order of
  Definition~\ref{def:dm-goals}.  The induction hypothesis is available for
  every \(A\)-, \(B\)-, or \(C\)-assertion whose multiset measure is smaller.
  The comparisons in that definition justify replacing a coordinate by an
  option, replacing one birthday by two smaller birthdays, and deleting an
  unused coordinate.

  Consider \(A(x,y)\).  By Definition~\ref{def:is-surreal}, it is enough to
  prove that every option of \(xy\) is surreal and that every left option is
  strictly below every right option.  Definition~\ref{def:game-mul} gives the
  four option families.  If \(x'\) and \(y'\) are options of \(x\) and \(y\),
  respectively, the corresponding option has the form
  \[
    \mulopt(x',y;x,y')=x'y+xy'-x'y'.
  \]
  The smaller assertions \(A(x',y)\), \(A(x,y')\), and \(A(x',y')\) prove
  that its three product terms are surreal.  Closure under addition and
  negation, Theorem~\ref{thm:surreal-add-closure}, then proves the whole
  expression surreal.

  For separation put \(P(u,v)=\mulopt(u,y;x,v)\).  By
  Definition~\ref{def:game-mul}, a left option of \(xy\) is
  \[
    P(x^L,y^L)\quad(LL),\qquad P(x^R,y^R)\quad(RR),
  \]
  and a right option is
  \[
    P(x^L,y^R)\quad(LR),\qquad P(x^R,y^L)\quad(RL).
  \]
  Consider first \(P(x^L_1,y^L)\) and \(P(x^L_2,y^R)\).  Since the first is a
  left option and the second is a right option, we need to prove
  \[
    P(x^L_1,y^L)\GameLt P(x^L_2,y^R).
  \]
  Theorem~\ref{thm:surreal-totality} compares \(x^L_1\) and \(x^L_2\).
  Also, Theorem~\ref{thm:between-options} and strict transitivity from
  Theorem~\ref{thm:game-order-calculus} give
  \(y^L\GameLt y\GameLt y^R\), hence \(y^L\GameLt y^R\).
  The three possible comparisons of the two left options give
  \[
    \begin{array}{ll}
    x^L_1\GameLt x^L_2:
      &P(x^L_1,y^L)\GameLt P(x^L_2,y^L)
         \GameLt P(x^L_2,y^R),\\[1mm]
    x^L_1\GameEq x^L_2:
      &P(x^L_1,y^L)\GameLt P(x^L_1,y^R)
         \GameLe P(x^L_2,y^R),\\[1mm]
    x^L_2\GameLt x^L_1:
      &P(x^L_1,y^L)\GameLt P(x^L_1,y^R)
         \GameLt P(x^L_2,y^R).
    \end{array}
  \]
  Here is the first branch in detail.  Assume \(x^L_1\GameLt x^L_2\).  The
  smaller assertion \(C(x^L_1,x^L_2,y)\) gives
  \[
    C_L(x^L_1,x^L_2;y,y^L)
      =\operatorname{Cross}_{\GameLt}(x^L_1,x^L_2;y^L,y).
  \]
  Substitute
  \((a,c;d,e;b)=(x^L_1,x^L_2;y^L,y;x)\) in rule \textup{(a)} of
  Theorem~\ref{thm:cross-calculus}.  It becomes
  \[
    \mulopt(x^L_1,y;x,y^L)\GameLt\mulopt(x^L_2,y;x,y^L),
    \qquad\text{that is,}\qquad
    P(x^L_1,y^L)\GameLt P(x^L_2,y^L).
  \]
  For the second step, the smaller assertion \(C(y^L,y^R,x)\), evaluated at
  the left option \(x^L_2\) of \(x\), gives
  \[
    \operatorname{Cross}_{\GameLt}(y^L,y^R;x^L_2,x).
  \]
  Applying rule \textup{(e)} of Theorem~\ref{thm:cross-calculus} gives
  \(\operatorname{Cross}_{\GameLt}(x^L_2,x;y^L,y^R)\).  Rule \textup{(c)},
  now with
  \((a,c;d,e;b)=(x^L_2,x;y^L,y^R;y)\), gives
  \[
    P(x^L_2,y^L)\GameLt P(x^L_2,y^R).
  \]
  Strict transitivity therefore proves the first chain in the display.

  If \(x^L_1\GameEq x^L_2\), the same \(y\)-slot argument first gives
  \(P(x^L_1,y^L)\GameLt P(x^L_1,y^R)\).  The two smaller assertions
  \(B(x^L_1,x^L_2,y)\) and \(B(x^L_1,x^L_2,y^R)\) give
  \[
    x^L_1y\GameEq x^L_2y,
    \qquad
    x^L_1y^R\GameEq x^L_2y^R.
  \]
  Hence, in \(\Games/{\GameEq}\),
  \[
    [P(x^L_1,y^R)]
      =[x^L_1y+xy^R-x^L_1y^R]
      =[x^L_2y+xy^R-x^L_2y^R]
      =[P(x^L_2,y^R)].
  \]
  Thus \(P(x^L_1,y^R)\GameEq P(x^L_2,y^R)\), and the strict--weak
  composition law proves the middle chain.

  Finally, assume \(x^L_2\GameLt x^L_1\).  After the same strict
  \(y\)-slot movement at \(x^L_1\), the smaller assertion
  \(C(x^L_2,x^L_1,y)\) gives
  \[
    C_R(x^L_2,x^L_1;y,y^R)
      =\operatorname{Cross}_{\GameLt}(x^L_2,x^L_1;y,y^R).
  \]
  Rule \textup{(b)} of Theorem~\ref{thm:cross-calculus},
  with \((a,c;d,e;b)=(x^L_2,x^L_1;y,y^R;x)\), reverses the \(x\)-slot
  direction and gives
  \[
    P(x^L_1,y^R)\GameLt P(x^L_2,y^R).
  \]
  Strict transitivity proves the last chain.  All the \(B\)- and
  \(C\)-assertions just invoked are smaller by the one-to-two birthday
  replacement in Definition~\ref{def:dm-goals}.

  The \(LL\)-versus-\(RL\), \(RR\)-versus-\(LR\), and
  \(RR\)-versus-\(RL\) comparisons are proved analogously.  This proves
  \(A(x,y)\).

  Now consider \(B(a,b,y)\).  Assume \(a\GameEq b\).  By
  Definition~\ref{def:game-order}, to prove \(ay\GameLe by\) one must rule out
  \(by\GameLe\lambda\) for every left option \(\lambda\) of \(ay\), and rule
  out \(\rho\GameLe ay\) for every right option \(\rho\) of \(by\).

  Definition~\ref{def:game-mul} gives two kinds of left option of \(ay\):
  \[
    \begin{aligned}
      \lambda_{LL}&=\mulopt(a^L,y;a,y^L)
        =a^Ly+ay^L-a^Ly^L,\\
      \lambda_{RR}&=\mulopt(a^R,y;a,y^R)
        =a^Ry+ay^R-a^Ry^R.
    \end{aligned}
  \]
  We must prove \(\lambda_{LL}\GameLt by\) and
  \(\lambda_{RR}\GameLt by\); either strict inequality contains the
  corresponding conclusion \(\neg(by\GameLe\lambda)\).

  For \(\lambda_{LL}\), the smaller assertion \(B(a,b,y^L)\) gives
  \[
    ay^L\GameEq by^L.
  \]
  The comparison \(a^L\GameLt a\), together with \(a\GameLe b\) from
  \(a\GameEq b\), gives \(a^L\GameLt b\).  Hence the smaller assertion
  \(C(a^L,b,y)\) gives the explicit cross inequality
  \[
    C_L(a^L,b;y,y^L):\qquad
    a^Ly+by^L\GameLt a^Ly^L+by.
  \]
  These are exactly the two hypotheses of the \(LL\) implication in
  Lemma~\ref{lem:mulopt-bounds}, so
  \[
    \lambda_{LL}=a^Ly+ay^L-a^Ly^L\GameLt by.
  \]

  For \(\lambda_{RR}\), the smaller assertion \(B(a,b,y^R)\) gives
  \[
    ay^R\GameEq by^R.
  \]
  Now \(b\GameLe a\) follows from \(a\GameEq b\), while
  \(a\GameLt a^R\); hence \(b\GameLt a^R\).  The smaller assertion
  \(C(b,a^R,y)\) therefore gives
  \[
    C_R(b,a^R;y,y^R):\qquad
    by^R+a^Ry\GameLt by+a^Ry^R.
  \]
  The \(RR\) implication in Lemma~\ref{lem:mulopt-bounds} now yields
  \[
    \lambda_{RR}=a^Ry+ay^R-a^Ry^R\GameLt by.
  \]
  Every assertion used here is smaller by Definition~\ref{def:dm-goals} and
  Lemma~\ref{lem:birthday-options}.

  The two right-option cases are analogous.  Explicitly, for the two right
  options of \(by\), one must prove
  \[
    \begin{aligned}
      ay&\GameLt \rho_{LR}
        =\mulopt(b^L,y;b,y^R)=b^Ly+by^R-b^Ly^R,\\
      ay&\GameLt \rho_{RL}
        =\mulopt(b^R,y;b,y^L)=b^Ry+by^L-b^Ry^L.
    \end{aligned}
  \]
  These follow analogously from the \(LR\) and \(RL\) implications of
  Lemma~\ref{lem:mulopt-bounds}.  The four exclusions prove \(ay\GameLe by\).

  Interchanging \(a\) and \(b\) gives \(by\GameLe ay\).  This does not invoke
  the induction hypothesis at the unchanged measure: the equality
  \[
    \mu_B(a,b,y)=\mu_B(b,a,y)
  \]
  from Definition~\ref{def:dm-goals} identifies the two collections of
  strictly smaller assertions, and symmetry of \(a\GameEq b\) permits the
  same one-sided proof to be reused.  The two weak inequalities are exactly
  \(ay\GameEq by\) by Definition~\ref{def:game-equivalence}.

  Finally consider \(C(x_1,x_2,y)\), with \(x_1\GameLt x_2\).  The assertion
  to be proved is that, for every left option \(y^L\) and right option \(y^R\)
  of \(y\),
  \[
    \begin{aligned}
      C_L(x_1,x_2;y,y^L):\quad
        &x_1y+x_2y^L\GameLt x_1y^L+x_2y,\\
      C_R(x_1,x_2;y,y^R):\quad
        &x_1y^R+x_2y\GameLt x_1y+x_2y^R.
    \end{aligned}
  \]
  Unfolding Definition~\ref{def:game-order} gives the following bridge
  alternative: either a right option \(x_1^R\) satisfies
  \(x_1^R\GameLe x_2\), or a left option \(x_2^L\) satisfies
  \(x_1\GameLe x_2^L\).  Indeed, if neither existed, the two universal clauses
  in that definition would give \(x_2\GameLe x_1\), contradicting
  \(x_1\GameLt x_2\).

  Suppose first that \(x_1^R\GameLe x_2\).  The smaller assertion
  \(A(x_1,y)\) makes \(x_1y\) surreal.  Its product options are those of
  Definition~\ref{def:game-mul}.  For each left option \(y^L\) of \(y\), the
  \(RL\) product option lies above \(x_1y\), so
  \[
    x_1y\GameLt x_1^Ry+x_1y^L-x_1^Ry^L.
  \]
  Work in the quotient of Definition~\ref{def:game-quotient}, whose ordered additive group
  structure is given by Theorem~\ref{thm:game-quotient-ordered-add-group}.  Adding
  \([x_1^Ry^L]\) gives
  \[
    [x_1y]+[x_1^Ry^L]<[x_1^Ry]+[x_1y^L].
  \]
  Reflection of strict order from the quotient gives
  \[
    \begin{aligned}
      C_L(x_1,x_1^R;y,y^L)
        &=\operatorname{Cross}_{\GameLt}(x_1,x_1^R;y^L,y),\\
        &\text{i.e.}\quad
          x_1y+x_1^Ry^L\GameLt x_1y^L+x_1^Ry.
    \end{aligned}
  \]
  Similarly, for a right option \(y^R\),
  the \(RR\) product option lies below \(x_1y\):
  \[
    x_1^Ry+x_1y^R-x_1^Ry^R\GameLt x_1y.
  \]
  Adding \([x_1^Ry^R]\) and reflecting strict order gives
  \[
    \begin{aligned}
      C_R(x_1,x_1^R;y,y^R)
        &=\operatorname{Cross}_{\GameLt}(x_1,x_1^R;y,y^R),\\
        &\text{i.e.}\quad
          x_1y^R+x_1^Ry\GameLt x_1y+x_1^Ry^R.
    \end{aligned}
  \]
  These are called the adjacent conditions because their two endpoints are
  \(x_1\) and its option \(x_1^R\).  The required conditions displayed above
  have endpoints \(x_1,x_2\), so it remains to move the second endpoint from
  \(x_1^R\) to \(x_2\).

  By Theorem~\ref{thm:surreal-totality}, either
  \(x_1^R\GameLt x_2\), \(x_1^R\GameEq x_2\), or
  \(x_2\GameLt x_1^R\).  If \(x_1^R\GameLt x_2\), the smaller assertion
  \(C(x_1^R,x_2,y)\) supplies
  \[
    \begin{aligned}
      C_L(x_1^R,x_2;y,y^L)
        &=\operatorname{Cross}_{\GameLt}(x_1^R,x_2;y^L,y),\\
        &\text{i.e.}\quad
          x_1^Ry+x_2y^L\GameLt x_1^Ry^L+x_2y,\\[1mm]
      C_R(x_1^R,x_2;y,y^R)
        &=\operatorname{Cross}_{\GameLt}(x_1^R,x_2;y,y^R),\\
        &\text{i.e.}\quad
          x_1^Ry^R+x_2y\GameLt x_1^Ry+x_2y^R.
    \end{aligned}
  \]
  Rule \textup{(f)} of Theorem~\ref{thm:cross-calculus} now composes
  \[
    \begin{gathered}
      \operatorname{Cross}_{\GameLt}(x_1,x_1^R;y^L,y)
      \quad\text{and}\quad
      \operatorname{Cross}_{\GameLt}(x_1^R,x_2;y^L,y),\\
      \operatorname{Cross}_{\GameLt}(x_1,x_1^R;y,y^R)
      \quad\text{and}\quad
      \operatorname{Cross}_{\GameLt}(x_1^R,x_2;y,y^R).
    \end{gathered}
  \]
  The results are precisely
  \(C_L(x_1,x_2;y,y^L)\) and \(C_R(x_1,x_2;y,y^R)\).

  If \(x_1^R\GameEq x_2\), the smaller \(B\)-assertions give
  \[
    x_1^Ry\GameEq x_2y,\qquad
    x_1^Ry^L\GameEq x_2y^L,\qquad
    x_1^Ry^R\GameEq x_2y^R.
  \]
  Replacing these three products in the two adjacent inequalities gives,
  respectively,
  \[
    x_1y+x_2y^L\GameLt x_1y^L+x_2y,
    \qquad
    x_1y^R+x_2y\GameLt x_1y+x_2y^R,
  \]
  which are again the two required conditions.  The remaining alternative
  \(x_2\GameLt x_1^R\) contradicts \(x_1^R\GameLe x_2\).

  The alternative \(x_1\GameLe x_2^L\) is analogous, using the adjacent
  conditions for \(x_2^L,x_2\) and then moving the first endpoint from
  \(x_2^L\) to \(x_1\).
\end{proof}

\begin{theorem}[Conway A: closure of the product]
  \label{thm:conway-a}
  \lean{Conway.conway_A}
  \leanok
  If \(x\) and \(y\) are surreal games, then \(xy\) is surreal.
\end{theorem}
\begin{proof}
  \leanok
  \uses{thm:abc-main}
  Apply the \(A(x,y)\) clause of Theorem~\ref{thm:abc-main} to the two
  surrealness hypotheses.  Its conclusion is precisely that \(xy\) is
  surreal.
\end{proof}

\begin{theorem}[Conway B: product respects equivalence]
  \label{thm:conway-b}
  \lean{Conway.conway_B}
  \leanok
  If \(x_1,x_2,y\) are surreal and \(x_1\GameEq x_2\), then
  \(x_1y\GameEq x_2y\).
\end{theorem}
\begin{proof}
  \leanok
  \uses{thm:abc-main}
  Apply the \(B(x_1,x_2,y)\) clause of Theorem~\ref{thm:abc-main}, with the
  three surrealness hypotheses and \(x_1\GameEq x_2\).  The conclusion is
  \(x_1y\GameEq x_2y\).
\end{proof}

\begin{theorem}[Conway C: strict cross inequalities]
  \label{thm:conway-c}
  \lean{Conway.conway_C}
  \leanok
  If \(x_1,x_2,y\) are surreal and \(x_1\GameLt x_2\), then all the conditions
  \(C_L(x_1,x_2;y,y^L)\) and \(C_R(x_1,x_2;y,y^R)\) hold.
\end{theorem}
\begin{proof}
  \leanok
  \uses{thm:abc-main}
  Apply the \(C(x_1,x_2,y)\) clause of Theorem~\ref{thm:abc-main}, with the
  three surrealness hypotheses and \(x_1\GameLt x_2\).  Its two conclusions
  are precisely the required left and right cross inequalities.
\end{proof}

\section{Congruence and associativity}

\begin{theorem}[Two-variable multiplicative congruence]
  \label{thm:game-mul-congr}
  \lean{Game.mul_congr}
  \leanok
  If \(x_1\GameEq x_2\) and \(y_1\GameEq y_2\), and all four games are surreal, then
  \[
    x_1y_1\GameEq x_2y_2.
  \]
\end{theorem}
\begin{proof}
  \leanok
  \uses{thm:game-order-calculus,thm:game-mul-laws,thm:conway-b}
  Conway \(B\), Theorem~\ref{thm:conway-b}, applied to
  \(x_1\GameEq x_2\) gives
  \(x_1y_1\GameEq x_2y_1\).  To replace the second factor, use product
  commutativity from Theorem~\ref{thm:game-mul-laws}, apply
  Theorem~\ref{thm:conway-b} to \(y_1\GameEq y_2\) with common factor \(x_2\),
  and commute back.  This gives
  \(x_2y_1\GameEq x_2y_2\).  Transitivity of game equivalence in
  Theorem~\ref{thm:game-order-calculus} composes the two replacements.
\end{proof}

\begin{theorem}[Associativity on surreal games]
  \label{thm:game-mul-assoc}
  \lean{Game.mul_assoc_of_isSurreal}
  \leanok
  If \(a,b,c\) are surreal games, then
  \[
    (ab)c\GameEq a(bc).
  \]
\end{theorem}
\begin{proof}
  \leanok
  \uses{def:game-inductions,lem:birthday-options,def:game-mul,def:game-neg,def:game-quotient,thm:game-quotient-ordered-add-group,def:is-surreal,thm:game-order-calculus,thm:option-extensionality,thm:game-neg-order,thm:game-mul-laws,thm:game-distrib,lem:game-mul-neg}
  Use the three-game well-founded induction of
  Definition~\ref{def:game-inductions}.  Lemma~\ref{lem:birthday-options}
  makes every option replacement smaller.  We strengthen the induction
  statement as follows: associativity may be used for any triple whose three
  birthdays are bounded by those of \((a,b,c)\), provided at least one bound
  is strict.  Thus, if \(a',b',c'\) are chosen options, the induction
  hypothesis contains all seven equivalences
  \[
    \begin{gathered}
      (a'b)c\GameEq a'(bc),\qquad
      (ab')c\GameEq a(b'c),\qquad
      (ab)c'\GameEq a(bc'),\\
      (a'b')c\GameEq a'(b'c),\qquad
      (a'b)c'\GameEq a'(bc'),\qquad
      (ab')c'\GameEq a(b'c'),\\
      (a'b')c'\GameEq a'(b'c').
    \end{gathered}
  \]

  We next compare a nested product option under the two bracketings.  For
  options \(a',b',c'\), we claim that:
  \[
    \mulopt\bigl(\mulopt(a',b;a,b'),c;ab,c'\bigr)
      \GameEq
    \mulopt\bigl(a',bc;a,\mulopt(b',c;b,c')\bigr).
  \]
  We prove this by expanding the left side and the right side separately.
  First, \(\mulopt(a',b;a,b')=a'b+ab'-a'b'\), so distributivity gives
  \[
    \begin{aligned}
      &\left[
        \mulopt\bigl(\mulopt(a',b;a,b'),c;ab,c'\bigr)
      \right]\\
      &\quad=
        \bigl[\mulopt(a',b;a,b')c\bigr]+[(ab)c']
          -\bigl[\mulopt(a',b;a,b')c'\bigr]\\
        &=[(a'b)c]+[(ab')c]-[(a'b')c]+[(ab)c']\\
        &\quad-[(a'b)c']-[(ab')c']+[(a'b')c'].
    \end{aligned}
  \]
  This is where the second equivalence of Lemma~\ref{lem:game-mul-neg} is
  used: right-distributivity produces the terms
  \[
    [(-(a'b'))c]
      =[-((a'b')c)],
    \qquad
    [(-(a'b'))c']
      =[-((a'b')c')].
  \]
  Substituting the seven strengthened induction hypotheses into this equality
  gives
  \[
    \begin{aligned}
      &\left[
        \mulopt\bigl(\mulopt(a',b;a,b'),c;ab,c'\bigr)
      \right]\\
        &\quad=[a'(bc)]+[a(b'c)]-[a'(b'c)]+[a(bc')]\\
        &\quad-[a'(bc')]-[a(b'c')]+[a'(b'c')].
    \end{aligned}
  \]

  For the other bracketing,
  \(\mulopt(b',c;b,c')=b'c+bc'-b'c'\).  Expanding directly gives
  \[
    \begin{aligned}
      &\left[
        \mulopt\bigl(a',bc;a,\mulopt(b',c;b,c')\bigr)
      \right]\\
      &\quad=[a'(bc)]
        +\bigl[a\mulopt(b',c;b,c')\bigr]
        -\bigl[a'\mulopt(b',c;b,c')\bigr]\\
        &=[a'(bc)]+[a(b'c)]+[a(bc')]-[a(b'c')]\\
        &\quad-[a'(b'c)]-[a'(bc')]+[a'(b'c')].
    \end{aligned}
  \]
  Here the first equivalence of Lemma~\ref{lem:game-mul-neg} is used in the
  two terms
  \[
    [a(-(b'c'))]=[-a(b'c')],
    \qquad
    [a'(-(b'c'))]=[-a'(b'c')].
  \]
  The right sides of the last two class expansions differ only by the order
  and association of their summands.  The additive commutative group structure
  of Theorem~\ref{thm:game-quotient-ordered-add-group} therefore gives
  \[
    \mulopt\bigl(\mulopt(a',b;a,b'),c;ab,c'\bigr)
      \GameEq
    \mulopt\bigl(a',bc;a,\mulopt(b',c;b,c')\bigr).
  \]
  This proves the claim.

  We now apply this claim to the option families of
  Definition~\ref{def:game-mul}.  First choose an \(LL\) option of \((ab)c\)
  whose selected left option of \(ab\) is itself \(LL\).  For
  \(a^L\in A^L,b^L\in B^L,c^L\in C^L\), it is
  \[
    L=\mulopt\bigl(\mulopt(a^L,b;a,b^L),c;ab,c^L\bigr).
  \]
  The game \(\mulopt(b^L,c;b,c^L)\) is an \(LL\) left option of \(bc\), so
  \[
    L'=\mulopt\bigl(a^L,bc;a,\mulopt(b^L,c;b,c^L)\bigr)
  \]
  is an \(LL\) left option of \(a(bc)\).  The preceding calculation with
  \((a',b',c')=(a^L,b^L,c^L)\) gives \(L\GameEq L'\).

  For a mixed example, choose an \(LR\) right option of \((ab)c\), again
  using an \(LL\) left option of \(ab\).  With \(c^R\in C^R\), it is
  \[
    R=\mulopt\bigl(\mulopt(a^L,b;a,b^L),c;ab,c^R\bigr).
  \]
  Now \(\mulopt(b^L,c;b,c^R)\) is an \(LR\) right option of \(bc\), and
  \[
    R'=\mulopt\bigl(a^L,bc;a,\mulopt(b^L,c;b,c^R)\bigr)
  \]
  is the corresponding \(LR\) right option of \(a(bc)\).  The calculation
  with \((a',b',c')=(a^L,b^L,c^R)\) gives \(R\GameEq R'\).

  The remaining cases choose \(a^R,b^R,c^L,c^R\) according to the outer and
  inner \(LL,RR,LR,RL\) families.  Reversing the calculation matches options
  from \(a(bc)\) back to \((ab)c\).  Definition~\ref{def:is-surreal} shows
  that every chosen option of \(a,b,c\) is surreal, so the seven strengthened
  recursive associativity statements apply.  Thus the left and right option
  lists match up to equivalence in both directions.
  Theorem~\ref{thm:option-extensionality} yields
  \((ab)c\GameEq a(bc)\).
\end{proof}


%%FILE blueprint/src/ordered_ring.tex
\chapter{The quotient ring and its order}

\section{Multiplication on equivalence classes}

\begin{theorem}[Multiplication on surreal representatives]
  \label{thm:surreal-mul}
  \lean{Surreal.mul_isSurreal,Surreal.mul,Surreal.mul_congr}
  \leanok
  If \(a,b\) are surreal representatives, their game product is surreal.  Moreover, replacing
  either representative by an equivalent one does not change the equivalence class of the
  product.  Thus multiplication is a well-defined operation on representatives modulo
  \(\GameEq\).
\end{theorem}
\begin{proof}
  \leanok
  \uses{def:surreal-representative,def:is-surreal,def:game-mul,def:game-order,def:game-equivalence,thm:conway-a,thm:game-mul-congr}
  By Definition~\ref{def:surreal-representative}, write
  \[
    a=(a_{\mathrm{game}},a_{\mathrm{proof}}),\qquad
    b=(b_{\mathrm{game}},b_{\mathrm{proof}}).
  \]
  Thus \(a_{\mathrm{proof}}\) and \(b_{\mathrm{proof}}\) witness that the short
  games \(a_{\mathrm{game}}\) and \(b_{\mathrm{game}}\) satisfy
  Definition~\ref{def:is-surreal}.  Theorem~\ref{thm:conway-a} shows that the
  product \(a_{\mathrm{game}}b_{\mathrm{game}}\), formed as in
  Definition~\ref{def:game-mul}, is again surreal.  Hence
  \[
    \bigl(a_{\mathrm{game}}b_{\mathrm{game}},(ab)_{\mathrm{proof}}\bigr)
  \]
  is a surreal representative, where \((ab)_{\mathrm{proof}}\) denotes the
  resulting witness.

  For alternative representatives
  \[
    a'=(a'_{\mathrm{game}},a'_{\mathrm{proof}}),\qquad
    b'=(b'_{\mathrm{game}},b'_{\mathrm{proof}}),
  \]
  assume that their underlying games satisfy
  \(a_{\mathrm{game}}\GameEq a'_{\mathrm{game}}\) and
  \(b_{\mathrm{game}}\GameEq b'_{\mathrm{game}}\) in the sense of
  Definition~\ref{def:game-equivalence}.  Theorem~\ref{thm:game-mul-congr}
  then gives
  \[
    a_{\mathrm{game}}b_{\mathrm{game}}
      \GameEq
    a'_{\mathrm{game}}b'_{\mathrm{game}}.
  \]
  Thus the equivalence class of the product is independent of both chosen
  representatives.
\end{proof}
\begin{definition}[Multiplication on short surreal numbers]
  \label{def:surreal-number-mul}
  \lean{Surreal.SurrealNumber.mul,Surreal.instMulSurrealNumber}
  \leanok
  \uses{def:surreal-number,thm:surreal-mul}
  Define multiplication on \(\NoShort\) by taking the game product of any two representatives
  and then passing to the quotient.
\end{definition}

\begin{theorem}[Multiplicative laws on the quotient]
  \label{thm:quotient-mul-laws}
  \lean{Surreal.SurrealNumber.mul_comm,Surreal.SurrealNumber.mul_assoc,Surreal.SurrealNumber.one_mul,Surreal.SurrealNumber.mul_one,Surreal.SurrealNumber.zero_mul,Surreal.SurrealNumber.mul_zero}
  \leanok
  Multiplication on \(\NoShort\) is associative and commutative, has identity \(1\), and has
  \(0\) as an absorbing element.
\end{theorem}
\begin{proof}
  \leanok
  \uses{def:surreal-number,def:surreal-number-mul,thm:game-mul-laws,thm:game-mul-assoc}
  By Definition~\ref{def:surreal-number}, choose surreal representatives
  \(a_0,b_0,c_0\) for three arbitrary classes.  Definition~\ref{def:surreal-number-mul}
  identifies their products with the classes of the corresponding game
  products.  The commutativity, unit, and zero equivalences of
  Theorem~\ref{thm:game-mul-laws} therefore give
  \[
    [a_0b_0]=[b_0a_0],\qquad [1\cdot a_0]=[a_0]=[a_0\cdot1],
    \qquad [0\cdot a_0]=[0]=[a_0\cdot0].
  \]
  Likewise, Theorem~\ref{thm:game-mul-assoc}, applied to the surreal games
  \(a_0,b_0,c_0\), gives
  \[
    [(a_0b_0)c_0]=[a_0(b_0c_0)].
  \]
  Equality of each displayed pair of classes follows precisely because the
  representative games are equivalent.  Since the representatives were
  arbitrary, all the asserted laws hold on \(\NoShort\).
\end{proof}

\begin{theorem}[Distributivity on the quotient]
  \label{thm:quotient-distrib}
  \lean{Surreal.SurrealNumber.left_distrib,Surreal.SurrealNumber.right_distrib}
  \leanok
  For \(a,b,c\in\NoShort\),
  \[
    a(b+c)=ab+ac,
    \qquad
    (a+b)c=ac+bc.
  \]
\end{theorem}
\begin{proof}
  \leanok
  \uses{def:surreal-number,thm:game-distrib,def:surreal-number-mul,thm:quotient-mul-laws}
  Choose surreal representatives \(a_0,b_0,c_0\) as permitted by
  Definition~\ref{def:surreal-number}.  By
  Definition~\ref{def:surreal-number-mul}, left distributivity reduces to the
  equivalence
  \[
    a_0(b_0+c_0)\GameEq a_0b_0+a_0c_0,
  \]
  which is Theorem~\ref{thm:game-distrib}.  Equivalent representatives define
  the same class, so this proves \(a(b+c)=ab+ac\).

  For right distributivity, use the commutativity in
  Theorem~\ref{thm:quotient-mul-laws} to write
  \[
    (a+b)c=c(a+b)=ca+cb=ac+bc,
  \]
  where the middle equality is the left distributive law just proved and the
  last equality again uses commutativity.
\end{proof}

\begin{theorem}[Commutative ring structure]
  \label{thm:comm-ring}
  \lean{Surreal.instCommRingSurrealNumber}
  \leanok
  The short surreal numbers \(\NoShort\) form a commutative ring.
\end{theorem}
\begin{proof}
  \leanok
  \uses{thm:ordered-add-group,def:surreal-number-mul,thm:quotient-mul-laws,thm:quotient-distrib}
  Theorem~\ref{thm:ordered-add-group} supplies the additive commutative group,
  including associativity, commutativity, the zero element, and additive
  inverses.  Definition~\ref{def:surreal-number-mul} supplies multiplication.
  Its associativity, commutativity, identity element, and zero laws are exactly
  Theorem~\ref{thm:quotient-mul-laws}, while both distributive axioms are
  Theorem~\ref{thm:quotient-distrib}.  These are precisely the axioms of a
  commutative ring.
\end{proof}

\section{Positive products and ordered-ring compatibility}

\begin{theorem}[A product of positive surreal games is positive]
  \label{thm:game-mul-pos}
  \lean{Surreal.Game.mul_pos_of_isSurreal}
  \leanok
  If \(x,y\) are surreal games and \(0\GameLt x\), \(0\GameLt y\), then
  \[
    0\GameLt xy.
  \]
\end{theorem}
\begin{proof}
  \leanok
  \uses{def:game,def:game-order,def:game-inductions,lem:birthday-options,def:is-surreal,def:cross-conditions,thm:game-order-calculus,thm:game-add-laws,thm:game-mul-laws,thm:conway-b,thm:conway-c}
  First show that \(0\GameLt y\) forces a left option \(y^L\) with
  \(0\GameLe y^L\).  If no left option had this property, the left clause in
  Definition~\ref{def:game-order} would give \(y\GameLe0\); the other clause
  is vacuous because \(0\) has no right options.  This contradicts the strict
  hypothesis.

  Proceed by the well-founded scheme of
  Definition~\ref{def:game-inductions}, using
  Lemma~\ref{lem:birthday-options}, and choose such a \(y^L\).  By
  Definition~\ref{def:is-surreal}, the option \(y^L\) is surreal.  Split according as
  \(y^L\GameLe0\).  If it does, then the two weak comparisons give
  \(y^L\GameEq0\).  Theorem~\ref{thm:conway-b}, together with the
  commutativity and zero-product laws in Theorem~\ref{thm:game-mul-laws}, gives
  \(xy^L\GameEq0\), hence \(0\GameLe xy^L\).  Otherwise
  Definition~\ref{def:game-order}, together with the already known comparison
  \(0\GameLe y^L\), gives \(0\GameLt y^L\).  Since \(y^L\) has smaller birthday, the induction
  hypothesis gives \(0\GameLt xy^L\), and therefore again
  \(0\GameLe xy^L\).

  Apply Theorem~\ref{thm:conway-c} to \(0\GameLt x\), with common surreal
  factor \(y\) and left option \(y^L\).  Simplifying its left cross inequality
  by the zero-product laws of Theorem~\ref{thm:game-mul-laws} and the zero-sum
  laws of Theorem~\ref{thm:game-add-laws} yields
  \(xy^L\GameLt xy\).  The mixed transitivity assertion of
  Theorem~\ref{thm:game-order-calculus}, applied to
  \(0\GameLe xy^L\GameLt xy\), proves \(0\GameLt xy\).
\end{proof}

\begin{theorem}[Positive and nonnegative products on the quotient]
  \label{thm:quotient-mul-pos}
  \lean{Surreal.SurrealNumber.zero_le_one,Surreal.SurrealNumber.mul_pos,Surreal.SurrealNumber.mul_nonneg,Surreal.SurrealNumber.instZeroLEOneClass}
  \leanok
  For \(a,b\in\NoShort\),
  \[
    0<a,\ 0<b\Longrightarrow 0<ab,
    \qquad
    0\leq a,\ 0\leq b\Longrightarrow 0\leq ab,
  \]
  and \(0\leq1\).
\end{theorem}
\begin{proof}
  \leanok
  \uses{def:game,def:game-order,def:surreal-number,thm:ordered-add-group,def:surreal-number-mul,thm:quotient-mul-laws,thm:game-mul-pos}
  For strict positivity, choose surreal representatives \(a_0,b_0\) of the
  two classes using Definition~\ref{def:surreal-number}.  The hypotheses become
  \(0\GameLt a_0\) and \(0\GameLt b_0\).  By
  Theorem~\ref{thm:game-mul-pos}, \(0\GameLt a_0b_0\); passing to equivalence
  classes and using Definition~\ref{def:surreal-number-mul} gives \(0<ab\).

  Suppose only \(0\leq a\) and \(0\leq b\).  The linear order of
  Theorem~\ref{thm:ordered-add-group} splits each weak comparison into a strict
  inequality or equality.  If both inequalities are strict, the first part of
  this theorem gives \(0<ab\), hence \(0\leq ab\).  If \(a=0\) or \(b=0\),
  the appropriate zero-product law in Theorem~\ref{thm:quotient-mul-laws}
  gives \(ab=0\).  These cases prove nonnegative-product closure.

  Finally, Definitions~\ref{def:game} and~\ref{def:game-order} show directly
  that \(0\GameLe1\): the left-option condition for \(0\) and the right-option
  condition for \(1\) are both vacuous.  Hence their equivalence classes
  satisfy \(0\leq1\).
\end{proof}

\begin{theorem}[Monotonicity of multiplication]
  \label{thm:quotient-mul-mono}
  \lean{Surreal.SurrealNumber.mul_le_mul_of_nonneg_left',Surreal.SurrealNumber.mul_le_mul_of_nonneg_right',Surreal.SurrealNumber.mul_lt_mul_of_pos_left',Surreal.SurrealNumber.mul_lt_mul_of_pos_right',Surreal.SurrealNumber.instPosMulMono,Surreal.SurrealNumber.instMulPosMono,Surreal.SurrealNumber.instPosMulStrictMono,Surreal.SurrealNumber.instMulPosStrictMono}
  \leanok
  Multiplication by a nonnegative short surreal number preserves weak inequalities on either
  side, and multiplication by a positive short surreal number preserves strict inequalities on
  either side.
\end{theorem}
\begin{proof}
  \leanok
  \uses{thm:ordered-add-group,thm:comm-ring,thm:quotient-mul-pos}
  Assume \(b\leq c\) and \(0\leq a\).  By
  Theorem~\ref{thm:ordered-add-group}, the first hypothesis is equivalent to
  \(0\leq c-b\).  The nonnegative-product assertion of
  Theorem~\ref{thm:quotient-mul-pos} gives \(0\leq a(c-b)\), while the ring
  laws of Theorem~\ref{thm:comm-ring} give \(a(c-b)=ac-ab\).  The ordered
  additive group equivalence
  \(0\leq ac-ab\Longleftrightarrow ab\leq ac\) yields weak monotonicity on
  the left.  For the right-sided assertion, suppose \(a\leq b\) and
  \(0\leq c\).  Then \(0\leq b-a\), so
  \[
    0\leq(b-a)c=bc-ac,
  \]
  and hence \(ac\leq bc\).

  If \(b<c\) and \(0<a\), use instead \(0<c-b\).  Strict positivity of a
  product, again by Theorem~\ref{thm:quotient-mul-pos}, gives
  \(0<a(c-b)=ac-ab\), which Theorem~\ref{thm:ordered-add-group} identifies
  with \(ab<ac\).  Similarly, \(a<b\) and \(0<c\) imply
  \[
    0<(b-a)c=bc-ac,
  \]
  and hence \(ac<bc\).  These four implications give the four weak and strict,
  left- and right-sided multiplication-order compatibilities asserted in the
  theorem.
\end{proof}

\begin{theorem}[Ordered- and strictly ordered-ring compatibility]
  \label{thm:ordered-ring}
  \lean{Surreal.instIsOrderedRingSurrealNumber,Surreal.instIsStrictOrderedRingSurrealNumber}
  \leanok
  The ring \(\NoShort\) is nontrivial; addition reflects order, multiplication by a
  nonnegative element preserves weak inequalities, and multiplication by a positive element
  preserves strict inequalities, on either side.
\end{theorem}
\begin{proof}
  \leanok
  \uses{def:game,def:game-order,def:surreal-number,thm:game-order-calculus,thm:ordered-add-group,thm:comm-ring,thm:quotient-mul-mono}
  Theorem~\ref{thm:ordered-add-group} gives translation invariance of the
  order, Theorem~\ref{thm:comm-ring} gives the commutative ring laws, and the
  weak and strict multiplication compatibility assertions are
  Theorem~\ref{thm:quotient-mul-mono}.  In particular, if
  \(c+a\leq c+b\), translating both sides by \(-c\) and simplifying with the
  additive group laws of Theorem~\ref{thm:ordered-add-group} gives
  \(a\leq b\).  Thus addition reflects order as required.

  It remains to supply nontriviality.  At game level, \(0\) is the sole left
  option of \(1\) by Definition~\ref{def:game}.  If \(1\GameLe0\), the left
  clause of Definition~\ref{def:game-order} would require
  \(\neg(0\GameLe0)\), contradicting the reflexivity assertion of
  Theorem~\ref{thm:game-order-calculus}.  Therefore \(0\) and \(1\) are not
  equivalent, so their classes in Definition~\ref{def:surreal-number} are
  distinct.  Order reflection and the weak and strict multiplication
  inequalities above establish all the claimed ordered-ring
  compatibility axioms.
\end{proof}

\begin{corollary}[Main endpoint: a linear ordered commutative ring]
  \label{thm:main-endpoint}
  \lean{Surreal.instLinearOrderSurrealNumber,Surreal.instCommRingSurrealNumber,Surreal.instIsStrictOrderedRingSurrealNumber}
  \leanok
  The short surreal numbers constructed above form a linear ordered commutative ring.
\end{corollary}
\begin{proof}
  \leanok
  \uses{thm:ordered-add-group,thm:comm-ring,thm:ordered-ring}
  By Theorem~\ref{thm:ordered-add-group}, \(\NoShort\) is a linear ordered
  additive commutative group.  By Theorem~\ref{thm:comm-ring}, \(\NoShort\) is
  a commutative ring.
  Theorem~\ref{thm:ordered-ring} proves compatibility of these algebraic
  operations with the linear order.  Hence \(\NoShort\) is a linear ordered
  commutative ring.
\end{proof}

